% !TeX root = ../../Book5.tex
% 纲要标题：## 第3章　整性、Noether 正规化与维数
% 纲要位置：几何学计划纲要.md，第 5111 行。
% 正文状态：初稿；尚未正式审核。

% 纲要内容开始
% 
% > 本章目标与证明范围：严格按整性工具、有限投影、维数比较的次序展开。普通正规化和 Noether 正规化的对象不同，章名相近不意味着同一过程。
% 
% 纲要内容结束

\chapter{整性、Noether 正规化与维数}
\label{g5:ch:03}
\BookChapterNumber{g5:ch:03}

\begin{chapterabstract}
	首一方程使一个新函数的高次幂能够用低次幂表示，因而把代数关系转化为模的有限性。本章从这个观察建立整扩张的基本性质，证明素理想的 lying-over、going-up、incomparability 及整闭基底上的 going-down 定理。随后在任意基域上证明 Noether 正规化，将有限型代数表示为多项式子代数上的有限模。素链定义的维数与代数无关参数的数目由此相接；主理想定理和广义高度定理进一步控制方程能够产生的高度。参数系、零维环与纤维计算把这些结论带回具体的局部模型。
\end{chapterabstract}

\begin{learninggoals}
	\begin{itemize}
		\item 用首一方程与有限模判据判断整性，证明其传递性，并区分整闭包与 Noether 正规化。
		\item 证明 lying-over、incomparability、going-up 和整闭基底上的 going-down，辨认各项假设的用途。
		\item 完成 Noether 正规化的消元证明，解释无限域上的线性代换及任意域上的幂次代换。
		\item 从素链定义高度和维数，证明主理想定理、广义高度定理、有限型整代数的维数公式及整扩张的维数保持。
		\item 构造局部参数系，比较维数、嵌入维数与有限长度，并计算有限映射及一般映射的纤维环。
	\end{itemize}
\end{learninggoals}

在平面曲线 \(y^2=x^3\) 上，给定 \(x\) 后，\(y\) 满足首一方程 \(T^2-x^3=0\)。因此坐标环中的每个函数都能写成 \(a(x)+b(x)y\)：两个生成方向就足以控制全部函数。这个现象既表达了投影到 \(x\) 轴的有限性，也提示了维数的来源。函数 \(x\) 可以自由变化，而 \(y\) 只在 \(x\) 上增加代数关系，所以曲线只有一个代数无关的参数。

\ProseParagraphBreak

换成双曲线 \(xy=1\)，投影到 \(x\) 轴时，每个非空纤维仍只有一个点，坐标环却是 \(k[x,x^{-1}]\)，不能由有限多个元素作为 \(k[x]\) 模生成。缺少 \(x=0\) 的纤维，正与无法用首一方程控制 \(x^{-1}\) 相呼应。若改取 \(t=x+y\)，则 \(x\) 满足 \(X^2-tX+1=0\)，有限性便恢复了。选择合适的投影方向，是 Noether 正规化要解决的问题。

\ProseParagraphBreak

本章先证明整扩张如何控制函数与素理想，再构造有限投影，最后比较参数数目与素链长度。第二章的局部化、有限模、行列式技巧和 Nakayama 引理将贯穿这些证明。一般代数论证允许任意基域；只有把闭点直接当作坐标取值时才另行假设基域代数闭。首一关系、素链以及局部切片分别记录不同的信息，后面的算例会把它们放在同一个环中比较。

% !TeX root = ../../Book5.tex
% 纲要标题：### 3.1 整扩张及有限性
% 纲要位置：几何学计划纲要.md，第 4821 行。
% 正文状态：未撰写；本文件目前是框架。

% 纲要内容开始
% 
% * 首一方程与有限模判据
% * 整性传递和整闭包
% * lying-over 与 incomparability
% * going-up 及整闭基底下的 going-down
% 
% 纲要内容结束

\section{整扩张及有限性}
\label{g5:sec:0301}

待撰写

\subsection{首一方程与有限模判据}
\label{g5:subsec:0301:01}
待撰写

\subsection{整性传递和整闭包}
\label{g5:subsec:0301:02}
待撰写

\subsection{lying-over 与 incomparability}
\label{g5:subsec:0301:03}
待撰写

\subsection{going-up 及整闭基底下的 going-down}
\label{g5:subsec:0301:04}
待撰写
% 纲要条目排为编号小节标题，并进入目录与书签；说明段仍保留于纲要注释中。

% !TeX root = ../../Book5.tex
% 纲要标题：### 3.2 Noether 正规化的证明
% 纲要位置：几何学计划纲要.md，第 4828 行。
% 正文状态：未撰写；本文件目前是框架。

% 纲要内容开始
% 
% * 代数依赖的变换消元
% * 无限基域的线性投影
% * 有限基域的非线性变换
% * 有限多项式子代数
% * 几何中的有限投影
% 
% 纲要内容结束

\section{Noether 正规化的证明}
\label{g5:sec:0302}

待撰写

\subsection{代数依赖的变换消元}
\label{g5:subsec:0302:01}
待撰写

\subsection{无限基域的线性投影}
\label{g5:subsec:0302:02}
待撰写

\subsection{有限基域的非线性变换}
\label{g5:subsec:0302:03}
待撰写

\subsection{有限多项式子代数}
\label{g5:subsec:0302:04}
待撰写

\subsection{几何中的有限投影}
\label{g5:subsec:0302:05}
待撰写
% 纲要条目排为编号小节标题，并进入目录与书签；说明段仍保留于纲要注释中。

% !TeX root = ../../Book5.tex
% 纲要标题：### 3.3 Krull 维数与高度
% 纲要位置：几何学计划纲要.md，第 5130 行。
% 正文状态：初稿；尚未正式审核。

% 纲要内容开始
% 
% * 素链、局部维数和高度
% * 主理想定理的证明
% * 广义高度定理
% * 有限型代数维数与超越次数
% * 整扩张的维数保持
% 
% 纲要内容结束

\section{Krull 维数与高度}
\label{g5:sec:0303}

一条不可约曲线可以包含点，一个不可约曲面又可以包含曲线。沿这样的严格包含逐次下降，最多能走多少步，是一种不依赖坐标选择的维数定义。在素谱中，不可约闭集写成 \(V(\mathfrak p)\)，而闭集的包含方向与素理想的包含方向相反。因此，维数问题可以转化为素理想链的长度问题。

\ProseParagraphBreak

上一节的正规化给出有限个代数无关的参数。本节将证明，这些参数的数目恰好等于素链所测出的维数。另一项核心问题是：用 \(r\) 个方程切出的不可约分支，能够具有多大的高度？回答它需要 \term[Krullweishu]{Krull 维数}[Krull dimension]及高度定理。\footnote{Wolfgang Krull（1899--1971），德国数学家，建立了交换环的维数理论，并对赋值论与理想理论作出重要贡献。} 维数的基本定义与有限型整代数的结论可对照\cite[第5节，第30--35页]{Matsumura1989}；高度定理的陈述可对照\cite[第13节，定理13.5]{Matsumura1989}。

\subsection{素链、局部维数和高度}
\label{g5:subsec:0303:01}

\begin{definition}\label{g5:def:krull-dimension}
	设 \(A\) 为交换含幺环。一条严格包含的有限\term[sulian]{素链}[chain of prime ideals]
	\[
	\mathfrak p_0\subsetneq\mathfrak p_1\subsetneq\cdots\subsetneq\mathfrak p_r
	\]
	的长度为 \(r\)，即严格包含的次数。所有素链长度的上确界称为 \(A\) 的 Krull 维数，记为 \(\dim A\)。对非零环，它属于 \(\{0,1,2,\ldots\}\cup\{\infty\}\)；约定零环的维数为 \(-\infty\)。
	对 \(\mathfrak p\in\operatorname{Spec}A\)，所有以 \(\mathfrak p\) 为末项的素链长度的上确界，称为 \(\mathfrak p\) 的\term[gaodu]{高度}[height]，记为 \(\operatorname{ht}_A\mathfrak p\)。对真理想 \(I\subsetneq A\)，再定义
	\[
	\operatorname{ht}_A I
	=\inf_{\mathfrak p\supseteq I}\operatorname{ht}_A\mathfrak p.
	\]
\end{definition}
\symidx{\(\dim A\)：环的 Krull 维数}
\symidx{\(\operatorname{ht}_A\mathfrak p\)：素理想的高度}

长度为零的链只有一个素理想。因此，非零环维数为零，恰表示它的素理想之间没有严格包含，也即每个素理想都是极大理想。对整环，这又等价于环本身为域：零理想是素理想，若它也是极大理想，商环就是原环。

\ProseParagraphBreak

这个定义也给出了 \(\operatorname{Spec}A\) 的拓扑维数。事实上，每个非空不可约闭集都是某个 \(V(\mathfrak p)\)。若 \(V(I)\) 不可约且 \(fg\in\sqrt I\)，则
\[
	V(I)\subseteq V(fg)=V(f)\cup V(g),
\]
不可约性使 \(V(I)\) 包含于其中一边，由\cref{g5:prop:radical-prime-intersection}便有 \(f\in\sqrt I\) 或 \(g\in\sqrt I\)，所以 \(\sqrt I\) 为素理想；反向已经在\secref{g5:subsec:0201:04}说明。再由根式与素理想的对应，素链恰好给出严格下降的非空不可约闭集链。

\begin{proposition}\label{g5:prop:dimension-localization}
	设 \(A\) 为非零交换含幺环，\(\mathfrak p\in\operatorname{Spec}A\)。则
	\begin{equation}\label{g5:eq:height-local-dimension}
	\begin{aligned}
	\operatorname{ht}_A\mathfrak p&=\dim A_{\mathfrak p},\\
	\dim A&=\sup_{\mathfrak q\in\operatorname{Spec}A}\operatorname{ht}_A\mathfrak q
	=\sup_{\mathfrak m\text{ 为极大理想}}\dim A_{\mathfrak m}.
	\end{aligned}
	\end{equation}
	若 \(I\subsetneq A\)，则 \(\dim(A/I)\leq\dim A\)。特别地，
	\[
	\dim\bigl(A/\sqrt{(0)}\bigr)=\dim A.
	\]
\end{proposition}
\begin{proof}
	由\cref{g5:thm:prime-localization}，\(A_{\mathfrak p}\) 的素理想与 \(A\) 中包含于 \(\mathfrak p\) 的素理想保序对应。\(A_{\mathfrak p}\) 是局部环，它的每条素链都能在末端接上极大理想 \(\mathfrak pA_{\mathfrak p}\)，所以其维数就是以该极大理想为末项的链长上确界。这给出第一个等式。\(A\) 中任意素链的末项都包含于某个极大理想，把这个极大理想接到链的末端便得到其余等式。
	商环 \(A/I\) 的素理想与包含 \(I\) 的素理想保序对应，故商环中的链也是原环中的链。每个素理想都包含 \(\sqrt{(0)}\)，所以商去幂零根时素链一条也没有删去，维数不变。
\end{proof}

\(\dim A_{\mathfrak p}\) 称为环在 \(\mathfrak p\) 处的\term[jubuweishu]{局部维数}[local dimension]。它计算的是到达 \(\mathfrak p\) 以前的素链；\(\dim(A/\mathfrak p)\) 则计算从 \(\mathfrak p\) 出发向上的素链，也就是闭集 \(V(\mathfrak p)\) 的维数。例如在 \(k[t]\) 中，零素理想处的局部环是 \(k(t)\)，局部维数为零，而其闭包是整个一维素谱。对 \(a\in k\)，闭点 \((t-a)\) 的高度为一，其闭包只有一个点，维数为零。

\ProseParagraphBreak

这里 \(\dim k[t]=1\) 只用了单变量多项式环是主理想整环：它的每个非零素理想都由不可约多项式生成，因此极大。对 \(k[x_1,\ldots,x_n]\)，目前先能写出长度为 \(n\) 的素链
\begin{equation}\label{g5:eq:coordinate-prime-chain}
	(0)\subsetneq(x_1)\subsetneq(x_1,x_2)
	\subsetneq\cdots\subsetneq(x_1,\ldots,x_n).
\end{equation}
每一步的商环仍是多项式整环，故这些理想确实为素理想。它给出维数的下界；证明没有更长的素链，则需要本节后面的定理。

\begin{proposition}\label{g5:prop:dimension-minimal-components}
	设 \(A\) 为非零 Noether 环，\(\mathfrak p_1,\ldots,\mathfrak p_s\) 为其全部极小素理想。则
	\[
	\dim A=\max_{1\leq i\leq s}\dim(A/\mathfrak p_i).
	\]
\end{proposition}
\begin{proof}
	\cref{g5:cor:minimal-support-associated}应用于有限模 \(A\)，说明极小素理想只有有限多个，而且每个素理想都包含其中一个。给定 \(A\) 中的任意素链，取包含于链首项的 \(\mathfrak p_i\)。整条链都包含 \(\mathfrak p_i\)，所以在 \(A/\mathfrak p_i\) 中给出同样长度的链。反之，每个商环中的素链都提升为 \(A\) 中的素链。对长度取上确界即得结论。
\end{proof}

\subsection{主理想定理的证明}
\label{g5:subsec:0303:02}

一个方程 \(f=0\) 的不可约分支对应于 \((f)\) 上的极小素理想。\term[zhulixiangdingli]{主理想定理}[principal ideal theorem]断言，在 Noether 环中，这些素理想的高度至多为一。证明的关键是先把所考察的分支局部化，使它成为唯一的闭点，再比较素理想的幂。下面先建立证明所需的下降链条件。

\begin{lemma}\label{g5:lem:local-radical-artinian-quotient}
	设 \((R,\mathfrak m)\) 为 Noether 局部环，\(J\subsetneq R\) 满足 \(\sqrt J=\mathfrak m\)。则存在 \(N\geq1\) 使 \(\mathfrak m^N\subseteq J\)，并且 \(R/J\) 的理想下降链稳定。
\end{lemma}
\begin{proof}
	取 \(\mathfrak m\) 的有限组生成元 \(u_1,\ldots,u_s\)。若 \(s=0\)，则 \(R\) 为域且 \(J=0\)，结论成立。否则，由 \(u_i\in\sqrt J\)，可取 \(e_i\geq1\) 使 \(u_i^{e_i}\in J\)。令
	\[
	N=1+\sum_{i=1}^s(e_i-1).
	\]
	任意总次数为 \(N\) 的 \(u_i\) 单项式都有一个指数至少为相应的 \(e_i\)，所以属于 \(J\)。这些单项式生成 \(\mathfrak m^N\)，故 \(\mathfrak m^N\subseteq J\)。

	令 \(\kappa=R/\mathfrak m\)，考察有限滤过
	\[
	F_0=R/J\supseteq F_1\supseteq\cdots\supseteq F_N=0,
	\qquad F_i=(\mathfrak m^i+J)/J.
	\]
	每个 \(F_i/F_{i+1}\) 都被 \(\mathfrak m\) 杀掉，又是有限 \(R\) 模，因而是有限维 \(\kappa\) 向量空间。给定 \(R/J\) 的理想下降链 \(L_1\supseteq L_2\supseteq\cdots\)，各
	\[
	\frac{(L_j\cap F_i)+F_{i+1}}{F_{i+1}}
	\subseteq F_i/F_{i+1}
	\]
	随 \(j\) 构成向量子空间下降链。有限维性及 \(i\) 的有限性保证这些链在同一个指标以后全部稳定。

	如果相邻两项 \(L'\supseteq L''\) 在每层具有相同的像，则它们相等。确实，由 \(F_N=0\) 开始向上归纳：若已知 \(L'\cap F_{i+1}=L''\cap F_{i+1}\)，对任意 \(v\in L'\cap F_i\)，同像性给出 \(w\in L''\cap F_i\) 使 \(v-w\in F_{i+1}\)，归纳假设又给出 \(v-w\in L''\)，于是 \(v\in L''\cap F_i\)。取 \(i=0\) 即得 \(L'=L''\)。因此原下降链也稳定。
\end{proof}

这个引理把一个环上的理想链，化成有限多个剩余域向量空间中的子空间链。以下证明会把它用于 \(J=(a)\)。为使模 \(a\) 后得到的关系能重新提升到原环，我们先把素理想的幂局部化再收缩；这样得到的理想对乘以 \(a\) 自动具有消去性质。

\begin{lemma}\label{g5:lem:principal-local-domain}
	设 \((R,\mathfrak m)\) 为 Noether 局部整环，\(a\in\mathfrak m\) 且 \(\sqrt{(a)}=\mathfrak m\)。则 \(\dim R\leq1\)。
\end{lemma}
\begin{proof}
	若 \(a=0\)，则整环中 \(\sqrt{(0)}=0\)，所以 \(\mathfrak m=0\)，\(R\) 为域。以下设 \(a\ne0\)。只需证明每个严格包含于 \(\mathfrak m\) 的素理想 \(\mathfrak p\) 都为零。这样的 \(\mathfrak p\) 不含 \(a\)，否则它包含 \(\sqrt{(a)}=\mathfrak m\)。对 \(n\geq1\)，令
	\[
	Q_n=\mathfrak p^nR_{\mathfrak p}\cap R.
	\]
	这些理想满足 \(Q_{n+1}\subseteq Q_n\)，而 \(a/1\) 在 \(R_{\mathfrak p}\) 中为单位，故
	\begin{equation}\label{g5:eq:symbolic-power-cancellation}
	ac\in Q_n\quad\Longrightarrow\quad c\in Q_n.
	\end{equation}
	由\cref{g5:lem:local-radical-artinian-quotient}，\(R/(a)\) 中的下降链
	\[
	\frac{Q_1+(a)}{(a)}\supseteq
	\frac{Q_2+(a)}{(a)}\supseteq\cdots
	\]
	稳定。因此存在 \(n\geq1\) 使 \(Q_n+(a)=Q_{n+1}+(a)\)。对 \(v\in Q_n\)，可写 \(v=w+ac\)，其中 \(w\in Q_{n+1}\)。于是 \(ac\in Q_n\)，由\eqref{g5:eq:symbolic-power-cancellation}得 \(c\in Q_n\)，所以
	\[
	Q_n=Q_{n+1}+aQ_n.
	\]

	商模 \(Q_n/Q_{n+1}\) 有限生成，因为 \(R\) Noether；它等于自身的 \(a\) 倍，从而等于自身的 \(\mathfrak m\) 倍。由\cref{g5:thm:nakayama-local}，这个商模为零，即 \(Q_n=Q_{n+1}\)。

	现在在 \(\mathfrak p\) 处局部化。局部化理想的收缩再扩张恢复原理想，故
	\[
	\mathfrak p^nR_{\mathfrak p}=Q_nR_{\mathfrak p}
	=Q_{n+1}R_{\mathfrak p}=\mathfrak p^{n+1}R_{\mathfrak p}.
	\]
	\(R_{\mathfrak p}\) 也是 Noether 局部环。对有限模 \(\mathfrak p^nR_{\mathfrak p}\) 再用一次 Nakayama 引理，得到 \(\mathfrak p^nR_{\mathfrak p}=0\)。若 \(b\in\mathfrak p\)，则 \((b/1)^n=0\)。局部化后的环仍是整环，而且 \(R\rightarrow R_{\mathfrak p}\) 单射，所以 \(b=0\)。这证明 \(\mathfrak p=0\)，从而素链长度至多为一。
\end{proof}

上述 \(Q_n\) 称为 \(\mathfrak p\) 的第 \(n\) 个\term[fuhaomi]{符号幂}[symbolic power]，通常记为 \(\mathfrak p^{(n)}\)。一般有 \(\mathfrak p^n\subseteq\mathfrak p^{(n)}\)，两者未必相等。这里需要符号幂，是因为 \(a\notin\mathfrak p\) 使 \(a\) 在局部化后成为单位，从而保证\eqref{g5:eq:symbolic-power-cancellation}。证明中的两次 Nakayama 引理分别消去 \(Q_n/Q_{n+1}\) 和 \(\mathfrak p^nR_{\mathfrak p}\)，作用的局部环也分别是 \(R\) 和 \(R_{\mathfrak p}\)。

\begin{theorem}[Krull 主理想定理]\label{g5:thm:principal-ideal}
	设 \(A\) 为 Noether 环，\(f\in A\)。若 \(\mathfrak p\) 是包含 \((f)\) 的极小素理想，则
	\[
	\operatorname{ht}_A\mathfrak p\leq1.
	\]
	特别地，若 \(A\) 为整环且 \(f\ne0\) 不是单位，则每个这样的 \(\mathfrak p\) 都有高度一。
\end{theorem}
\begin{proof}
	令 \(R=A_{\mathfrak p}\)、\(\mathfrak m=\mathfrak pA_{\mathfrak p}\)。局部化的素理想对应说明，\(\mathfrak m\) 是 \(R\) 中唯一包含 \(fR\) 的素理想，所以 \(\sqrt{fR}=\mathfrak m\)。若 \(\dim R\geq2\)，便存在素链
	\[
	\mathfrak q_0\subsetneq\mathfrak q_1\subsetneq\mathfrak m.
	\]
	商环 \(R/\mathfrak q_0\) 是 Noether 局部整环；\(f\) 的像 \(\overline f\) 所生成理想的根式等于 \(\mathfrak m/\mathfrak q_0\)，因为 \(fR\subseteq fR+\mathfrak q_0\subseteq\mathfrak m\)。\cref{g5:lem:principal-local-domain}给出 \(\dim(R/\mathfrak q_0)\leq1\)，却又有长度为二的素链
	\[
	(0)\subsetneq\mathfrak q_1/\mathfrak q_0
	\subsetneq\mathfrak m/\mathfrak q_0,
	\]
	矛盾。因此 \(\operatorname{ht}_A\mathfrak p=\dim R\leq1\)。在整环且 \(f\ne0\) 的情形，\(\mathfrak p\ne0\)，素链 \((0)\subsetneq\mathfrak p\) 又给出高度至少为一。
\end{proof}

极小性不可省去。例如 \(k[x,y]\) 中的 \((x,y)\) 包含主理想 \((x)\)，但它下方存在素链 \((0)\subsetneq(x)\subsetneq(x,y)\)，高度至少为二。主理想定理控制的是方程 \(x=0\) 的整个不可约分支，而非分支上每一个更小闭子集的高度。

\subsection{广义高度定理}
\label{g5:subsec:0303:03}

一个方程所产生的高度界，怎样推广到多个方程？关键是把任意一条素链变成商环中的素链，并控制能够损失的长度。下面的引理完成这一点；它也将说明局部环中逐次加入方程时，维数每次至多下降一。

\begin{lemma}\label{g5:lem:principal-dimension-drop}
	设 \((R,\mathfrak m)\) 为 Noether 局部环，\(a\in\mathfrak m\)。则
	\[
	\dim R\leq\dim\bigl(R/(a)\bigr)+1.
	\]
\end{lemma}
\begin{proof}
	给定一条以极大理想为末项的素链
	\[
	\mathfrak q_0\subsetneq\mathfrak q_1\subsetneq\cdots
	\subsetneq\mathfrak q_d=\mathfrak m.
	\]
	我们将构造一条全部包含 \(a\)、长度至少为 \(d-1\) 的素链。它在 \(R/(a)\) 中仍严格，便得到所需估计。

	由于 \(a\in\mathfrak m\)，存在最小指标 \(t\) 使 \(a\in\mathfrak q_t\)。如果 \(t=0\)，原链的每一项都包含 \(a\)，结论已经成立。以下设 \(t\geq1\)。令 \(\mathfrak p_t=\mathfrak q_t\)，再从 \(i=t-1\) 向 \(i=0\) 递归选择
	\[
	\mathfrak q_i+(a)\subseteq\mathfrak p_i
	\subseteq\mathfrak p_{i+1},
	\]
	其中 \(\mathfrak p_i\) 为 \(\mathfrak q_i+(a)\) 上的一个极小素理想。这样的选择总能完成：\(\mathfrak p_{i+1}\) 包含 \(\mathfrak q_i+(a)\)，而 Noether 商环 \(R/(\mathfrak q_i+(a))\) 中的每个素理想都包含一个极小素理想，见\cref{g5:cor:minimal-support-associated}。

	对 \(0\leq i\leq t-2\)，有 \(\mathfrak p_i\subsetneq\mathfrak p_{i+1}\)。若相等，因为 \(a\notin\mathfrak q_{i+1}\)，就有
	\[
	\mathfrak q_i\subsetneq\mathfrak q_{i+1}
	\subsetneq\mathfrak p_i.
	\]
	但 \(\mathfrak p_i/\mathfrak q_i\) 是 Noether 环 \(R/\mathfrak q_i\) 中主理想 \((\overline a)\) 上的极小素理想，\cref{g5:thm:principal-ideal}要求它的高度至多为一，与上面的长度二素链矛盾。

	因此
	\[
	\mathfrak p_0\subsetneq\cdots\subsetneq\mathfrak p_{t-1}
	\subseteq\mathfrak q_t\subsetneq\cdots\subsetneq\mathfrak q_d
	\]
	全部包含 \(a\)，且只有中间的 \(\mathfrak p_{t-1}\subseteq\mathfrak q_t\) 可能成为等号。删去可能的重复项后，所得素链的长度至少为 \(d-1\)。于是每个这样的 \(d\) 都满足
	\(d\leq\dim(R/(a))+1\)。任意素链都能接到 \(\mathfrak m\)，对链长取上确界即得结论；这个论证也涵盖维数为无穷的可能情形。
\end{proof}

证明中向下选择 \(\mathfrak p_i\)，是为了始终保留素理想之间的包含关系。主理想定理随后保证：原链在遇到包含 \(a\) 的项以前，相邻两层不会被合并。把这个逐链估计在商环中反复使用，就得到\term[Krullgaodudingli]{Krull 高度定理}[Krull's height theorem]。

\begin{theorem}[Krull 高度定理]\label{g5:thm:height-theorem}
	设 \(A\) 为 Noether 环，真理想 \(I=(f_1,\ldots,f_r)\) 由 \(r\) 个元素生成。若 \(\mathfrak p\) 是 \(I\) 上的极小素理想，则
	\[
	\operatorname{ht}_A\mathfrak p\leq r.
	\]
	因而 \(\operatorname{ht}_A I\leq r\)。允许 \(r=0\)，此时 \(I=(0)\)。
\end{theorem}
\begin{proof}
	在 \(\mathfrak p\) 处局部化，记 \(R=A_{\mathfrak p}\)、\(\mathfrak m=\mathfrak pA_{\mathfrak p}\)。由于 \(\mathfrak p\) 在 \(I\) 上极小，\(\mathfrak m\) 是 \(R\) 中唯一包含 \(IR\) 的素理想。因此 \(R/IR\) 的素谱只有一个点，维数为零。

	依次记
	\[
	R_j=R/(f_1,\ldots,f_j)R\qquad(0\leq j\leq r),
	\]
	其中 \(R_0=R\)。各 \(R_j\) 都是非零 Noether 局部环，因为这些生成元都属于 \(\mathfrak m\)。对 \(R_{j-1}\) 中 \(f_j\) 的像应用\cref{g5:lem:principal-dimension-drop}，得到
	\[
	\dim R_{j-1}\leq\dim R_j+1.
	\]
	从 \(\dim R_r=0\) 开始向前归纳，便有 \(\dim R\leq r\)。由\eqref{g5:eq:height-local-dimension}，这正是所求的高度界。\(r=0\) 时，\(R=R/IR\) 已为零维，同样成立。
\end{proof}

这个结论也称\term[guangyigaodudingli]{广义高度定理}[generalized principal ideal theorem]。主理想定理对应于 \(r=1\)，而 \(k[x_1,\ldots,x_r]\) 中的 \((x_1,\ldots,x_r)\) 表明上界可以取到：\eqref{g5:eq:coordinate-prime-chain}中的坐标素链给出高度至少为 \(r\)，高度定理给出高度至多为 \(r\)。

\begin{corollary}\label{g5:cor:noetherian-local-finite-dimension}
	设 \((R,\mathfrak m)\) 为 Noether 局部环，且 \(\mathfrak m\) 能由 \(r\) 个元素生成。则 \(\dim R\leq r<\infty\)。特别地，Noether 环的每个素理想都有有限高度。
\end{corollary}
\begin{proof}
	\(\mathfrak m\) 是自身上唯一的极小素理想，故高度定理给出
	\(\dim R=\operatorname{ht}_R\mathfrak m\leq r\)。对一般 Noether 环的素理想 \(\mathfrak p\)，局部化 \(A_{\mathfrak p}\) 仍为 Noether 环，其极大理想有限生成，再用同一结论即可。
\end{proof}

维数因此给出了方程数目的下界：若 \((R,\mathfrak m)\) 为维数 \(d\) 的 Noether 局部环，且 \(\sqrt{(f_1,\ldots,f_r)}=\mathfrak m\)，则 \(r\geq d\)。下一节将构造恰好 \(d\) 个这样的元素，并把它们称为参数系。

\subsection{有限型代数维数与超越次数}
\label{g5:subsec:0303:04}

上一节把有限型代数表示为多项式子代数上的有限模。对于整环，其参数数目由分式域的超越次数确定。为了把这个数与素链长度联系起来，先说明严格增加一个素理想会使超越次数严格下降。下面的论证适用于任意基域，所用的只是代数独立与有限生成。

\begin{lemma}\label{g5:lem:strict-prime-transcendence-drop}
	设 \(k\) 为域，\(B\) 为有限型 \(k\) 代数且为整环，\(\mathfrak q\ne(0)\) 为 \(B\) 的素理想。则
	\[
	\operatorname{trdeg}_k\operatorname{Frac}(B/\mathfrak q)
	<\operatorname{trdeg}_k\operatorname{Frac}(B).
	\]
\end{lemma}
\begin{proof}
	写 \(B=k[b_1,\ldots,b_s]\)。从 \(b_i\) 在 \(B/\mathfrak q\) 中的像里选取一个极大代数无关子集，重排指标后记为 \(\overline b_1,\ldots,\overline b_e\)。其余生成元都在 \(k(\overline b_1,\ldots,\overline b_e)\) 上代数，所以这些像构成 \(\operatorname{Frac}(B/\mathfrak q)\) 的一个超越基。\(b_1,\ldots,b_e\) 本身也代数无关：原环中的多项式关系在商环中仍成立。因此
	\[
	e\leq d=\operatorname{trdeg}_k\operatorname{Frac}(B).
	\]
	若 \(e=d\)，则 \(\operatorname{Frac}(B)\) 在 \(K=k(b_1,\ldots,b_e)\) 上代数；否则再加入一个超越于 \(K\) 的元素，就得到 \(d+1\) 个代数无关元，与\cref{g5:lem:transcendence-bound}矛盾。令
	\[
	S=k[b_1,\ldots,b_e]\setminus\{0\}.
	\]
	因为 \(\overline b_1,\ldots,\overline b_e\) 代数无关，\(S\cap\mathfrak q=\varnothing\)。另一方面，
	\[
	S^{-1}B=K[b_{e+1},\ldots,b_s]
	\]
	是由有限多个代数元生成的 \(K\) 代数，故是有限维 \(K\) 向量空间：对各生成元取代数方程并用其约去高次幂，就得到有限个单项式张成。它又是整环，所以是域；事实上，任意非零元素的乘法为有限维向量空间的单射，因而满射。
	但 \(S\cap\mathfrak q=\varnothing\) 使 \(S^{-1}\mathfrak q\) 为这个域的真理想，而 \(\mathfrak q\ne0\) 及整环的局部化单射性又使 \(S^{-1}\mathfrak q\ne0\)，矛盾。因此 \(e<d\)。
\end{proof}

该证明的要点是：若取素商后没有损失代数无关参数，就能同时把这些参数中的非零多项式变为单位。原环此时变成有限代数域扩张，原来非零的素理想却必须仍为真理想，两者无法同时成立。

\begin{theorem}[有限型代数的维数]\label{g5:thm:affine-dimension}
	设 \(k\) 为任意域，\(A\) 为非零有限型 \(k\) 代数。
	若 \(A\) 为整环，则
	\begin{equation}\label{g5:eq:affine-dimension-transcendence}
	\dim A=\operatorname{trdeg}_k\operatorname{Frac}(A).
	\end{equation}
	一般地，若 \(\mathfrak p_1,\ldots,\mathfrak p_s\) 为 \(A\) 的全部极小素理想，则
	\begin{equation}\label{g5:eq:affine-dimension-components}
	\dim A=\max_{1\leq i\leq s}
	\operatorname{trdeg}_k\operatorname{Frac}(A/\mathfrak p_i).
	\end{equation}
\end{theorem}
\begin{proof}
	先设 \(A\) 为整环，并记 \(d=\operatorname{trdeg}_k\operatorname{Frac}(A)\)。对任意素链
	\[
	\mathfrak q_0\subsetneq\mathfrak q_1\subsetneq\cdots\subsetneq\mathfrak q_r,
	\]
	各 \(A/\mathfrak q_i\) 都是有限型 \(k\) 代数，且都是整环。对相邻两项应用\cref{g5:lem:strict-prime-transcendence-drop}，得到一列严格下降的非负整数
	\[
	\operatorname{trdeg}_k\operatorname{Frac}(A/\mathfrak q_0)
	>\cdots>
	\operatorname{trdeg}_k\operatorname{Frac}(A/\mathfrak q_r)\geq0.
	\]
	首项至多为 \(d\)：若 \(\mathfrak q_0=0\) 则相等，若 \(\mathfrak q_0\ne0\) 则再用一次该引理。故 \(r\leq d\)，从而 \(\dim A\leq d\)。

	反向用\cref{g5:thm:noether-normalization,g5:cor:normalization-transcendence}：存在代数无关的 \(t_1,\ldots,t_d\in A\)，使
	\[
	C=k[t_1,\ldots,t_d]\subseteq A
	\]
	为有限整扩张。\(C\) 中有长度为 \(d\) 的坐标素链
	\[
	(0)\subsetneq(t_1)\subsetneq\cdots\subsetneq(t_1,\ldots,t_d).
	\]
	当 \(d=0\) 时，这条链只含 \((0)\) 一项。由\cref{g5:thm:lying-over}，可先在 \(A\) 中选取位于 \((0)\) 上方的素理想；再反复应用\cref{g5:thm:going-up}，便把整条链提升到 \(A\)。各提升理想的收缩互不相同，所以提升后的包含也都严格，长度仍为 \(d\)。因此 \(\dim A\geq d\)，证明\eqref{g5:eq:affine-dimension-transcendence}。

	一般的 \(A\) 由\cref{g5:thm:hilbert-basis,g5:prop:affine-noetherian}知为 Noether 环。\cref{g5:prop:dimension-minimal-components}给出
	\(\dim A=\max_i\dim(A/\mathfrak p_i)\)，再对每个整环 \(A/\mathfrak p_i\) 使用已证的有限型代数结论，就得到\eqref{g5:eq:affine-dimension-components}。
\end{proof}

\begin{corollary}\label{g5:cor:polynomial-dimension}
	对任意域 \(k\) 和整数 \(n\geq0\)，有
	\[
	\dim k[x_1,\ldots,x_n]=n.
	\]
\end{corollary}
\begin{proof}
	分式域 \(k(x_1,\ldots,x_n)\) 对 \(k\) 的超越次数为 \(n\)，应用\cref{g5:thm:affine-dimension}即可。
\end{proof}

代数闭域 \(k\) 上的非空仿射代数集 \(X\)，其维数定义为 \(\dim X=\dim k[X]\)。坐标环同构保持素链，因此这个定义与仿射实现无关。若 \(X\) 不可约，就得到 \(\dim X=\operatorname{trdeg}_k k(X)\)；一般情形则取各不可约分量维数的最大值。这样，坐标中独立参数的数目与不可约闭集链的长度给出了同一个维数。

\begin{example}\label{g5:exm:mixed-dimensional-components}
	设 \(k\) 为任意域，令
	\[
	A=k[x,y,z]/(xy,xz).
	\]
	在 \(k[x,y,z]\) 中有
	\[
	(xy,xz)=(x)\cap(y,z).
	\]
	右端包含左端。反之，若 \(xg\in(y,z)\)，模去 \((y,z)\) 得到 \(x\overline g=0\)；商环 \(k[x]\) 是整环，故 \(\overline g=0\)，即 \(g\in(y,z)\)。这证明交理想等式。两个素理想 \((x)\)、\((y,z)\) 互不包含，而且任何包含其交的素理想都包含其中一个，所以它们恰是交理想上的极小素理想。
	相应两个分支的商环分别为
	\[
	A/\bigl((x)/(xy,xz)\bigr)\cong k[y,z],
	\qquad
	A/\bigl((y,z)/(xy,xz)\bigr)\cong k[x].
	\]
	由\cref{g5:thm:affine-dimension}，两者维数分别为二和一，而 \(\dim A=2\)。当 \(k\) 为代数闭域时，相应代数集为平面 \(x=0\) 与 \(x\) 轴，交于原点。
	令 \(\mathfrak p=(y,z)/(xy,xz)\)。它是极小素理想，所以 \(\operatorname{ht}_A\mathfrak p=0\)，但 \(\dim(A/\mathfrak p)=1\)。于是
	\[
	\operatorname{ht}_A\mathfrak p+\dim(A/\mathfrak p)=1<2=\dim A.
	\]
	这个计算说明，整体维数取决于最高维的分支，某个素理想的高度则只计算到达该素理想的链。两者相减一般不能给出其闭包的维数。
\end{example}

\subsection{整扩张的维数保持}
\label{g5:subsec:0303:05}

有限型代数的维数证明已经用到：多项式子环中的素链可以沿整扩张提升。对一般整扩张，lying-over 与 going-up 保证底环中的链不会丢失，incomparability 保证上层的严格链收缩后仍然严格。这两点合起来给出维数相等。

\begin{theorem}\label{g5:thm:integral-dimension}
	设 \(A\subseteq B\) 为非零交换含幺环的整扩张。则
	\[
	\dim B=\dim A,
	\]
	其中允许维数为无穷。对任意 \(\mathfrak q\in\operatorname{Spec}B\)，令 \(\mathfrak p=\mathfrak q\cap A\)，还有
	\[
	\dim(B/\mathfrak q)=\dim(A/\mathfrak p).
	\]
	更一般地，对任意交换含幺环同态 \(\varphi:A\rightarrow B\)，若 \(B\) 在 \(\varphi(A)\) 上整，则
	\[
	\dim B=\dim(A/\ker\varphi).
	\]
\end{theorem}
\begin{proof}
	设 \(\mathfrak q_0\subsetneq\cdots\subsetneq\mathfrak q_r\) 为 \(B\) 中的素链。收缩得到 \(A\) 中的包含链。若相邻两项收缩相同，就与\cref{g5:thm:incomparability}矛盾。因此收缩后的每一步都严格，长度仍为 \(r\)，从而 \(\dim B\leq\dim A\)。

	反过来，给定 \(A\) 中任意素链 \(\mathfrak p_0\subsetneq\cdots\subsetneq\mathfrak p_r\)，先由\cref{g5:thm:lying-over}选取 \(\mathfrak q_0\in\operatorname{Spec}B\)，使 \(\mathfrak q_0\cap A=\mathfrak p_0\)。再由\cref{g5:thm:going-up}逐步得到
	\[
	\mathfrak q_0\subseteq\mathfrak q_1\subseteq\cdots\subseteq\mathfrak q_r,
	\qquad \mathfrak q_i\cap A=\mathfrak p_i.
	\]
	收缩互不相同，故这些包含全部严格。因此 \(\dim B\geq\dim A\)。对所有有限链的长度取上确界，两个不等式也涵盖无穷维情形。

	给定 \(\mathfrak q\) 后，诱导包含 \(A/\mathfrak p\subseteq B/\mathfrak q\) 仍为整扩张，对它使用已经证明的结论即得第二个等式。对于一般整同态，先用 \(A/\ker\varphi\cong\varphi(A)\) 代替底环，再对 \(\varphi(A)\subseteq B\) 使用结论。若 \(B=0\)，两边均为零环的维数 \(-\infty\)。
\end{proof}

这个维数等式适用于任意整扩张。整性通过素理想的上行与不可比性直接控制全部有限素链。若希望比较指定点处的高度，则还需要保留提升链的上端点；这正是 going-down 所处理的问题。

\begin{proposition}\label{g5:prop:integral-height-comparison}
	设 \(A\subseteq B\) 为非零交换含幺环的整扩张，\(\mathfrak q\in\operatorname{Spec}B\)，\(\mathfrak p=\mathfrak q\cap A\)。则
	\[
	\operatorname{ht}_B\mathfrak q\leq\operatorname{ht}_A\mathfrak p.
	\]
	若这个扩张还满足 going-down，则等号成立。特别地，若 \(A\) 是整闭整环、\(B\) 是整环，则等号成立。
\end{proposition}
\begin{proof}
	每条以 \(\mathfrak q\) 为末项的素链收缩后，依不可比性仍严格，且末项为 \(\mathfrak p\)，因此得到第一个不等式。若有 going-down，给定任意以 \(\mathfrak p\) 为末项的链
	\[
	\mathfrak p_0\subsetneq\cdots\subsetneq\mathfrak p_r=\mathfrak p,
	\]
	从指定的 \(\mathfrak q_r=\mathfrak q\) 开始，逐次向下选取 \(\mathfrak q_i\subseteq\mathfrak q_{i+1}\)，使其收缩为 \(\mathfrak p_i\)。收缩的严格性保证所选链也严格，于是 \(\operatorname{ht}_B\mathfrak q\geq r\)。取上确界即得反向不等式。最后的充分条件由\cref{g5:thm:going-down}给出。
\end{proof}

比较这两个证明可以看出端点的作用。证明整体维数相等时，先随意选定链首素理想的一个上方素理想，再向上提升；证明指定 \(\mathfrak q\) 的高度相等时，必须从这个已经给定的末项向下构造。虽然得到的都是同样长度的链，所需要的提升方向却不同。

\begin{example}\label{g5:exm:integral-dimension-cusp}
	设 \(k\) 为任意域。包含
	\[
	k[t^2,t^3]\subseteq k[t]
	\]
	为有限整扩张：\(t\) 满足首一方程 \(T^2-t^2=0\)，而 \(k[t]\) 由 \(1,t\) 作为 \(k[t^2,t^3]\) 模生成。因此
	\[
	\dim k[t^2,t^3]=\dim k[t]=1.
	\]
	把 \(u,v\) 分别映到 \(t^2,t^3\)，可得
	\[
	k[u,v]/(v^2-u^3)\cong k[t^2,t^3].
	\]
	确实，模去 \(v^2-u^3\) 后每个多项式可写为 \(a(u)+v b(u)\)；代入后，\(a(t^2)\) 只含偶次幂，\(t^3b(t^2)\) 只含奇次幂，两者之和为零便迫使 \(a=b=0\)。这既核对了核理想，也说明尖点曲线与其参数直线具有同一维数。
\end{example}

一般整同态中的核确实影响维数。例如满同态 \(k[x,y]\rightarrow k[y]\)、\(x\mapsto0\)、\(y\mapsto y\) 为整同态，每个目标元素都满足一个首一的一次方程。但源环维数为二、目标环维数为一，维数相等发生在
\[
	k[x,y]/(x)\cong k[y]
\]
与目标环之间。因而应用整扩张的维数保持时，先确定底环在目标环中的像，就确定了实际参与比较的空间。
% 纲要条目排为编号小节标题，并进入目录与书签；说明段仍保留于纲要注释中。

% !TeX root = ../../Book5.tex
% 纲要标题：### 3.4 维数的局部模型
% 纲要位置：几何学计划纲要.md，第 5138 行。
% 正文状态：初稿；尚未正式审核。

% 纲要内容开始
% 
% * 参数系及嵌入维数
% * 零维 Noether 环和有限长度
% * 超曲面与完全交的预告
% * 纤维维数问题的代数准备
% 
% ---
% 
% 纲要内容结束

\section{维数的局部模型}
\label{g5:sec:0304}

一条曲线在一个点附近需要几个参数，可能与描述这个点附近的一阶变化所需的坐标数不同。例如尖点的函数域只有一个超越方向，但原点处的两个坐标都不能从极大理想的极小生成元组中删去。维数、生成元数和方程数因此需要分别计算。本节先用参数系测量局部维数，再用长度记录零维部分的厚度，最后把这些量放进超曲面、完全交与变纤维的具体模型中。

\subsection{参数系及嵌入维数}
\label{g5:subsec:0304:01}

设 \((R,\mathfrak m,\kappa)\) 为非零 Noether 局部环。\(\mathfrak m\) 有限生成，由广义高度定理，它的高度 \(d=\dim R\) 有限。寻找参数的目标是选出尽量少的 \(x_i\in\mathfrak m\)，使商环 \(R/(x_1,\ldots,x_r)\) 的素谱只剩下闭点。这个条件等价于 \(\sqrt{(x_1,\ldots,x_r)}=\mathfrak m\)，因而允许商环仍有幂零元。

\begin{definition}
\label{g5:def:system-of-parameters}
	设 \((R,\mathfrak m,\kappa)\) 为非零 Noether 局部环，\(d=\dim R\)。若 \(x_1,\ldots,x_d\in\mathfrak m\) 满足
	\[
		\sqrt{(x_1,\ldots,x_d)}=\mathfrak m,
	\]
	就称它们为 \(R\) 的一组\term[canshuxi]{参数系}[system of parameters]。当 \(d=0\) 时允许空参数系。数
	\[
		\operatorname{edim}R\coloneqq
		\dim_\kappa(\mathfrak m/\mathfrak m^2)
	\]
	称为 \(R\) 的\term[qianruweishu]{嵌入维数}[embedding dimension]。
	\symidx{\(\operatorname{edim}R\)：局部环的嵌入维数}
\end{definition}

Nakayama 引理说明，嵌入维数是生成极大理想本身所需的最少元素数。参数系则只要求生成理想的根式等于极大理想。证明参数系存在时，我们需要选一个元素，同时避开有限多个不可约分支所对应的素理想。下面的\term[youxiansulixiangbikaiyinli]{有限素理想避开引理}[finite prime avoidance lemma]提供这一选择。

\begin{lemma}[有限素理想避开引理]
\label{g5:lem:finite-prime-avoidance}
	设 \(A\) 为交换含幺环，\(I\subseteq A\) 为理想，\(\mathfrak p_1,\ldots,\mathfrak p_s\) 为有限多个素理想。若 \(I\) 不包含在任何一个 \(\mathfrak p_i\) 中，则存在
	\[
		a\in I\setminus\bigcup_{i=1}^s\mathfrak p_i.
	\]
\end{lemma}
\begin{proof}
	删去包含在其他 \(\mathfrak p_j\) 中的素理想，不改变它们的并，所以可设所余素理想两两互不包含。对每个 \(i\)，选 \(a_i\in I\setminus\mathfrak p_i\)；对每个 \(j\ne i\)，选 \(b_{ij}\in\mathfrak p_j\setminus\mathfrak p_i\)。令
	\[
		c_i=a_i\prod_{j\ne i}b_{ij}.
	\]
	则 \(c_i\in I\)，且 \(c_i\notin\mathfrak p_i\)，因为 \(\mathfrak p_i\) 为素理想；对 \(j\ne i\)，又有 \(c_i\in\mathfrak p_j\)。所以 \(a=\sum_i c_i\) 在模 \(\mathfrak p_i\) 后等于非零的 \(c_i\)，从而避开全部 \(\mathfrak p_i\)。空集合情形任取 \(a\in I\) 即可。
\end{proof}

\begin{theorem}
\label{g5:thm:system-of-parameters}
	设 \((R,\mathfrak m)\) 为非零 Noether 局部环，\(d=\dim R\)。则参数系存在，并且
	\[
		d=\min\{r:\text{存在 }x_1,\ldots,x_r\in\mathfrak m,\quad
		\sqrt{(x_1,\ldots,x_r)}=\mathfrak m\}.
	\]
\end{theorem}
\begin{proof}
	先证明下界。若 \(r\) 个元素生成的理想根式为 \(\mathfrak m\)，则 \(\mathfrak m\) 是该理想的唯一极小素理想。\cref{g5:thm:height-theorem}给出
	\(\dim R=\operatorname{ht}\mathfrak m\leq r\)。

	下面对 \(d\) 归纳构造至多 \(d\) 个这样的元素。若 \(d=0\)，局部环的每个素理想都等于 \(\mathfrak m\)，故由根式等于所有包含它的素理想之交，得到 \(\sqrt{(0)}=\mathfrak m\)，空参数系满足要求。

	设 \(d>0\)。Noether 环只有有限多个极小素理想；这也可直接从\cref{g5:thm:prime-filtration}用于 \(R\) 时的有限支集并公式看出：每个极小素理想必须等于其中某个滤过素理想。取其中满足
	\(\dim R/\mathfrak p_i=d\)
	的那些。这样的素理想确实存在：取一条长度为 \(d\) 的素链，其最小端必为极小素理想，否则可在前面再添一项，违背 \(d\) 的定义。每个所选 \(\mathfrak p_i\) 都严格包含于 \(\mathfrak m\)，因为 \(\dim R/\mathfrak m=0\)。

	由\cref{g5:lem:finite-prime-avoidance}，可选 \(x\in\mathfrak m\) 避开这些 \(\mathfrak p_i\)。我们断言 \(\dim R/(x)\leq d-1\)。否则有一条长度为 \(d\) 的素链，其每项都包含 \(x\)。这条链的最小端 \(\mathfrak q\) 同样必为极小素理想，且 \(\dim R/\mathfrak q=d\)，于是 \(\mathfrak q\) 属于被避开的那一组；但 \(x\in\mathfrak q\)，矛盾。

	商环 \(R/(x)\) 仍为非零 Noether 局部环，设其维数为 \(d'\leq d-1\)。由归纳假设，在商环中可以选 \(d'\) 个元素生成以 \(\mathfrak m/(x)\) 为根式的理想。将它们提升回 \(R\)，连同 \(x\) 一共得到 \(1+d'\leq d\) 个元素，其生成理想的根式为 \(\mathfrak m\)。前面的下界又要求元素数至少为 \(d\)，故恰好为 \(d\)，得到参数系。
\end{proof}

参数系把素链定义的维数转换成了“切到只剩一个点所需的最少方程数”。这里保留商环的幂零结构，因此“只剩一个点”与“商环是域”并不相同。上述结论可与\cite[第13节，定理13.4--13.6；第14节开头]{Matsumura1989}对照；这里由已经证明的高度定理得到下界，存在性则使用有限素理想避开和归纳。

\begin{proposition}
\label{g5:prop:dimension-embedding-bound}
	设 \((R,\mathfrak m,\kappa)\) 为非零 Noether 局部环。则
	\[
		\dim R\leq\operatorname{edim}R<\infty.
	\]
\end{proposition}
\begin{proof}
	\(\mathfrak m\) 是有限生成模，所以 \(\mathfrak m/\mathfrak m^2\) 有有限维数 \(e\)。提升其一组基，由\cref{g5:cor:nakayama-generators}得到 \(\mathfrak m\) 的 \(e\) 个生成元。它们生成的理想根式当然为 \(\mathfrak m\)，故\cref{g5:thm:system-of-parameters}给出 \(\dim R\leq e\)。
\end{proof}

例如，设 \(k\) 为域，
\[
	R=\bigl(k[x,y]/(y^2-x^3)\bigr)_{(x,y)}.
\]
由\cref{g5:exm:cusp-birational}，这个环的分式域是 \(k(t)\)，其中 \(x=t^2,y=t^3\)。仿射维数定理给出整条尖点曲线维数为 \(1\)；局部环中有严格素链 \((0)\subsetneq(x,y)R\)，所以 \(\dim R=1\)。另一方面，关系 \(y^2-x^3\) 没有一次项，\(x,y\) 的像在 \(\mathfrak m/\mathfrak m^2\) 中线性无关：若有线性关系，清去常数项非零的局部分母后比较一次项，就会使两个系数都为零。因此 \(\operatorname{edim}R=2\)。单个 \(x\) 已是一组参数，因为 \(y^2=x^3\) 使 \(y\in\sqrt{(x)}\)。参数数目与极大理想生成元数的差别，在这里由同一个方程明确呈现。

\subsection{零维 Noether 环和有限长度}
\label{g5:subsec:0304:02}

参数切片虽然只剩一个点，仍可能保留不同厚度。例如 \(k[t]/(t^2)\) 与 \(k[t]/(t^3)\) 都只有一个素理想，维数都是零，但它们作为 \(k\) 向量空间的维数分别为 \(2,3\)。在不预先指定基域的一般环上，长度以简单模为单位记录这种差别。

\begin{definition}
\label{g5:def:module-length}
	设 \(A\) 为交换含幺环，\(M\) 为 \(A\) 模。所有有限严格子模链
	\[
		0=M_0\subsetneq M_1\subsetneq\cdots\subsetneq M_r=M
	\]
	的长度 \(r\) 的上确界，记为 \(\ell_A(M)\)，称为 \(M\) 的长度。若这个上确界有限，就称 \(M\) 为有限长度模；约定 \(\ell_A(0)=0\)。非零且没有非零真子模的模称为\term[jiandanmo]{简单模}[simple module]。
	\symidx{\(\ell_A(M)\)：模的长度}
\end{definition}

这延续\subsecref{g5:subsec:0104:05}中有限商代数的长度。有限长度模中，一条达到最大长度的链不能再插入子模，所以它的每个相邻商都是简单模；这种链称为\term[zuchenglie]{组成列}[composition series]。反过来，存在有限组成列会给出有限长度，这可由下面的加性逐段证明。

\begin{proposition}
\label{g5:prop:length-additivity}
	设 \(A\) 为交换含幺环，且 \(0\rightarrow L\rightarrow M\rightarrow N\rightarrow0\) 为短正合列。则 \(M\) 有限长度，当且仅当 \(L,N\) 都有限长度；在此情形，
	\[
		\ell_A(M)=\ell_A(L)+\ell_A(N).
	\]
\end{proposition}
\begin{proof}
	将 \(L\) 看成 \(M\) 的子模。对 \(U\subseteq V\subseteq M\)，若 \(U\cap L=V\cap L\)，且 \(U,V\) 在 \(N\) 中的像相同，则 \(U=V\)：对 \(v\in V\)，取 \(u\in U\) 具有相同的像，便有 \(v-u\in V\cap L=U\cap L\)。

	因此，\(M\) 的严格子模链每前进一步，或者与 \(L\) 的交严格增大，或者在 \(N\) 中的像严格增大。如果 \(L,N\) 长度有限，链长至多为 \(\ell_A(L)+\ell_A(N)\)，所以 \(M\) 长度有限。反过来，\(L\) 的子模链可看成 \(M\) 中的链，\(N\) 的链可拉回 \(M\)，所以 \(M\) 长度有限时二者都有限。最后，将 \(L\) 中的一条最长链与 \(N\) 中一条最长链的逆像拼接，得到长度恰为两者之和的链，证明等式。
\end{proof}

简单 \(A\) 模都是 \(A/\mathfrak m\) 的形式，其中 \(\mathfrak m\) 极大：任取一个非零元素，它生成整个简单模，核便是极大理想。因此，在局部环上，每个组成因子都同构于剩余域。若局部环还是剩余域为 \(k\) 的有限维 \(k\) 代数，则每个组成因子恰有一维，故长度等于 \(k\) 维数。长度有限也立即使子模升链和降链都终止。加性与链条件的关系可对照\cite[第3节，定理3.1及其后的讨论]{Matsumura1989}。

\ProseParagraphBreak

满足理想降链条件的环称为\term[Artinhuan]{Artin 环}[Artinian ring]。\footnote{Emil Artin（1898--1962），奥地利裔美国数学家，研究域论、代数数论与具有降链条件的环。}Noether 条件控制升链；对交换环，零维性恰好补足降链条件。

\begin{theorem}
\label{g5:thm:artinian-zero-dimension}
	设 \(A\ne0\) 为交换含幺环。以下条件等价：
	\begin{equivalences}
		\item \(A\) 为 Noether 环且 \(\dim A=0\)；
		\item \(A\) 为 Artin 环；
		\item \(A\) 作为自身的模具有有限长度。
	\end{equivalences}
	在这些条件下，\(\operatorname{Spec}A\) 是有限集，所有素理想均极大。零环同样为 Noether 环、Artin 环且长度为零，其空素谱按前文的空集约定处理。
\end{theorem}
\begin{proof}
	先设 \(A\) Noether 且维数为零。对 \(A\) 使用\cref{g5:thm:prime-filtration}，得到有限个相邻商 \(A/\mathfrak p_i\)。每个 \(\mathfrak p_i\) 都极大，所以这些商都是简单模。长度加性使 \(\ell_A(A)\) 有限。有限长度又立即蕴含理想降链条件。

	现在设 \(A\) Artin。对任意素理想 \(\mathfrak p\)，整环 \(D=A/\mathfrak p\) 仍满足降链条件。若 \(0\ne a\in D\)，链 \((a)\supseteq(a^2)\supseteq\cdots\) 终止，故某个 \(n\) 满足 \(a^n=a^{n+1}b\)。在整环中约去 \(a^n\) 得 \(ab=1\)。所以 \(D\) 是域，全部素理想均极大。

	极大理想也只能有有限多个。否则可取两两不同的 \(\mathfrak m_1,\mathfrak m_2,\ldots\)。链
	\(\mathfrak m_1\supseteq\mathfrak m_1\cap\mathfrak m_2\supseteq\cdots\)
	每一步都严格：对第 \(r+1\) 个极大理想，从每个 \(\mathfrak m_i\setminus\mathfrak m_{r+1}\) 中取元素，乘积属于前 \(r\) 个的交，却不属于素理想 \(\mathfrak m_{r+1}\)。这违背降链条件。记全部极大理想为 \(\mathfrak m_1,\ldots,\mathfrak m_s\)，并记
	\(\mathfrak n=\bigcap_i\mathfrak m_i=\sqrt{(0)}\)。

	我们证明 \(\mathfrak n\) 幂零，而不预先假设它有限生成。降链条件使某个 \(q\geq1\) 满足 \(\mathfrak n^q=\mathfrak n^{q+1}\)。令 \(I=\mathfrak n^q\)，则 \(I^2=I\)。若 \(I\ne0\)，在满足 \(IJ\ne0\) 的理想 \(J\) 中选一个极小者。存在 \(x\in J\) 使 \(Ix\ne0\)，故极小性使 \(J=Ax\)。另一方面，\(I(IJ)=IJ\ne0\)，极小性还给出 \(IJ=J\)。于是有限生成模 \(J\) 满足 \(IJ=J\)。由\cref{g5:prop:nakayama-ideal}，存在 \(a\in I\subseteq\mathfrak n\) 使 \((1-a)J=0\)。\(a\) 幂零，故 \(1-a\) 可逆，迫使 \(J=0\)，矛盾。因此 \(I=0\)。

	不同极大理想两两互素，中国剩余定理给出
	\[
		A/\mathfrak n\cong\prod_{i=1}^s A/\mathfrak m_i.
	\]
	每层 \(\mathfrak n^j/\mathfrak n^{j+1}\) 都是这个有限个域的乘积环上的模，并且继承降链条件。用乘积环的坐标幂等元分解它，得到有限个域上的向量空间。每个向量空间都必须有限维：无限基可以不断删去一个基向量，构造严格的下降子空间链。因此各层都是有限长度 \(A\) 模。由于 \(\mathfrak n^q=0\)，有限层的长度加性给出 \(\ell_A(A)<\infty\)。

	最后，有限长度使理想升链终止，所以 \(A\) Noether；又由刚才对 Artin 环素商的论证得 \(\dim A=0\)。全部等价关系得证。
\end{proof}

Artin 环推出有限长度的证明可对照\cite[第3节，定理3.2]{Matsumura1989}。在 Noether 局部环 \((R,\mathfrak m)\) 中，对真理想 \(J\) 因而有
\[
	\sqrt J=\mathfrak m
	\quad\Longleftrightarrow\quad
	\dim R/J=0
	\quad\Longleftrightarrow\quad
	\ell_R(R/J)<\infty.
\]
第一个等价来自局部环中的素理想对应，第二个来自上个定理；商环的子模就是它的理想，因此作为 \(R\) 模与作为 \(R/J\) 模计算的长度相同。特别地，参数系切出的局部商环总有有限长度。

\begin{example}
\label{g5:exm:parameter-lengths}
	设 \(k\) 为域，\(R=(k[x,y]/(y^2-x^3))_{(x,y)}\)。参数 \(x\) 与参数 \(y\) 分别给出
	\[
		R/(x)\cong k[y]/(y^2),\qquad
		R/(y)\cong k[x]/(x^3).
	\]
	这里局部化没有再改变右侧的环，因为它们已经是仅有一个极大理想的局部环。两个商的长度分别为 \(2,3\)。所以参数系的元素数只由维数决定，而参数切片的长度还取决于具体选择。

	另取任意域扩张 \(K/k\)。域 \(K\) 作为自身的模长度为 \(1\)，但 \(\dim_k K\) 可以无限。这说明有限长度不等于对任意外部基域都有限维；第一章有限点模型中的等同还使用了有限型及剩余域条件。
\end{example}

\subsection{超曲面与完全交的预告}
\label{g5:subsec:0304:03}

一个非零方程通常被期待削去一个维数。高度定理给出了一方面的控制：由 \(r\) 个元素生成的理想，其极小素理想高度至多为 \(r\)。要得到精确维数，还须核对方程所处的环及方程之间的关系。多项式环中的单个非恒定方程可以通过首一化得到一个完整模型。

\begin{proposition}
\label{g5:prop:polynomial-hypersurface-dimension}
	设 \(k\) 为任意域，\(n\geq1\)，\(f\in k[x_1,\ldots,x_n]\) 为非恒定多项式。则
	\[
		\dim k[x_1,\ldots,x_n]/(f)=n-1.
	\]
	更具体地，经过多项式坐标的可逆变换及乘以非零常数，可以使 \(f\) 关于最后一个变量首一；此时商环是其余 \(n-1\) 个变量组成的多项式子环上的有限自由模。
\end{proposition}
\begin{proof}
	由\cref{g5:lem:weighted-monic-change}，可以作一个对任意基域都有效的三角代换，使最后一个变量 \(T\) 的最高次系数为非零常数。除以该常数，写成
	\[
		f=T^e+a_1(Y)T^{e-1}+\cdots+a_e(Y),
		\qquad e\geq1,
	\]
	其中 \(Y=(Y_1,\ldots,Y_{n-1})\)。令 \(C=k[Y_1,\ldots,Y_{n-1}]\)。关于首一多项式作除法，每个剩余类都有唯一表达
	\[
		b_0(Y)+b_1(Y)T+\cdots+b_{e-1}(Y)T^{e-1}.
	\]
	唯一性来自非零首一多项式的倍式不可能有小于 \(e\) 的 \(T\) 次数。因此 \(C\) 单射到商环，且 \(1,T,\ldots,T^{e-1}\) 为其自由基。特别地，这是有限整扩张。由\cref{g5:thm:integral-dimension}及多项式环维数，商环维数为 \(n-1\)。
\end{proof}

在代数闭域上，由单个非恒定多项式定义的 \(V(f)\subseteq\mathbb A_k^n\) 称为\term[chaoqumian]{超曲面}[hypersurface]。取根式不会改变素谱的包含关系，所以 \(V(f)\) 的约化坐标环也有维数 \(n-1\)。商环 \(k[x_1,\ldots,x_n]/(f)\) 还保存方程的重数；例如 \(f=x_1^a\) 与 \(f=x_1\) 给出相同的点集，却一般给出不同的商环。

\ProseParagraphBreak

对多个方程，必须要求新方程在前面方程切出的商环中仍有独立作用。设 \(S\) 为交换含幺环，\(f_1,\ldots,f_r\in S\) 生成真理想。若对每个 \(i\)，乘以 \(f_i\) 在
\(S/(f_1,\ldots,f_{i-1})\)
上是单射，则称它们为一个\term[zhengzexulie]{正则序列}[regular sequence]。在给定的多项式环中，由正则序列生成的理想所定义的商代数是\term[wanquanjiao]{完全交}[complete intersection]的基本形式。后续将从局部方程研究其内在判别；这里先计算一类可以逐项验证的模型。

\begin{example}
\label{g5:exm:coordinate-complete-intersection}
	设 \(k\) 为域，\(S=k[x_1,\ldots,x_n]\)，\(0\leq r\leq n\)，并取整数 \(a_1,\ldots,a_r\geq1\)。序列
	\[
		x_1^{a_1},\ldots,x_r^{a_r}
	\]
	是正则序列。事实上，商去前 \(i-1\) 项后，所得环作为 \(k[x_i,\ldots,x_n]\) 模，有自由基
	\[
		x_1^{b_1}\cdots x_{i-1}^{b_{i-1}},
		\qquad 0\leq b_j<a_j.
	\]
	在这个自由模上乘以 \(x_i^{a_i}\) 单射，因为系数环是整环。最终商环
	\[
		B=S/(x_1^{a_1},\ldots,x_r^{a_r})
	\]
	是 \(k[x_{r+1},\ldots,x_n]\) 上的有限自由模，基为所有 \(0\leq b_i<a_i\) 的单项式 \(\prod_{i=1}^r x_i^{b_i}\)。该多项式子环单射到 \(B\)，所以整扩张保持维数给出
	\[
		\dim B=n-r.
	\]
	当 \(r=n\) 时，\(B\) 只支撑在原点，且
	\(\ell_B(B)=\dim_k B=\prod_i a_i\)。
	于是余维由被控制的坐标方向数决定，长度则还记住各方向的幂次。
\end{example}

两个最小生成方程也可能未使每个分支的余维都等于二。在代数闭域 \(k\) 上取
\[
	I=(xy,xz)\subseteq S=k[x,y,z].
\]
有
\[
	I=(x)\cap(y,z):
\]
若 \(xg\in(y,z)\)，模去 \((y,z)\) 后得到 \(x\overline g=0\) 于整环 \(k[x]\)，故 \(g\in(y,z)\)，证明了反向包含。因此零点集是平面 \(x=0\) 与直线 \(y=z=0\) 的并，两个分支的维数分别为 \(2,1\)，整个商环维数为 \(2\)。在原点的局部环中，\(xy,xz\) 仍是这个理想的极小生成元：其像在 \(I/\mathfrak mI\) 中线性无关，因为 \(\mathfrak mI\) 中没有非零二次齐次项。方程个数为二，并非因为多写了一个可删去的生成元。

\ProseParagraphBreak

这个序列失败在第二步。\(xy\) 在整环 \(S\) 中是非零因子，但 \(xz\) 在 \(S/(xy)\) 中被非零类 \(y\) 杀掉：
\(y(xz)=z(xy)=0\)。
新方程在某个已有分支上恒等于零，因而不能把那个分支再切掉一个维数。高度定理控制的是最小可能的方程数与素理想高度；具体商的维数还需要检查这些消失关系。

\subsection{纤维维数问题的代数准备}
\label{g5:subsec:0304:04}

在映射 \(A\rightarrow B\) 中，把基点固定在 \(\mathfrak p\in\operatorname{Spec}A\)，应当使 \(\mathfrak p\) 中的元素变为零，并使 \(A\setminus\mathfrak p\) 中的元素可逆。两个要求合起来给出\term[xianweihuan]{纤维环}[fibre ring]
\[
	B(\mathfrak p)\coloneqq B\otimes_A\kappa(\mathfrak p)
	\cong (A\setminus\mathfrak p)^{-1}(B/\mathfrak pB).
\]

\symidx{\(B(\mathfrak p)=B\otimes_A\kappa(\mathfrak p)\)：素点处的纤维环}
其素理想恰对应 \(B\) 中收缩为 \(\mathfrak p\) 的素理想：在商环中要求包含 \(\mathfrak pB\)，在局部化中要求避开 \(A\setminus\mathfrak p\)，所以收缩既包含又不大于 \(\mathfrak p\)。这只是商环及局部化的素理想对应，适用于任意环同态，纤维环也可能为零环。

\ProseParagraphBreak

若 \(A\) 是整环，零素理想的剩余域为 \(K=\operatorname{Frac}A\)。相应的 \(B\otimes_AK\) 称为\term[fanxianweihuan]{泛纤维环}[generic fibre ring]。它把所有非零基底函数都反演，因而适合测量映射在基底的泛点处还剩多少独立方向。纤维环与泛纤维的记法可对照\cite[第15节，Fibres，第116页]{Matsumura1989}。

\begin{proposition}
\label{g5:prop:generic-fiber-dimension}
	设 \(A\subseteq B\) 为整环的单射，且 \(B\) 是有限型 \(A\) 代数。令 \(K=\operatorname{Frac}A\)、\(L=\operatorname{Frac}B\)。则
	\[
		\dim(B\otimes_AK)=\operatorname{trdeg}_K L.
	\]
	若 \(A,B\) 还是某个域 \(k\) 上的有限型代数，则
	\[
		\dim(B\otimes_AK)=\dim B-\dim A.
	\]
\end{proposition}
\begin{proof}
	因为 \(A\subseteq B\) 单射且二者为整环，
	\[
		B\otimes_AK=(A\setminus\{0\})^{-1}B
	\]
	是非零有限型 \(K\) 整代数，其分式域仍为 \(L\)。对这个 \(K\) 代数应用\cref{g5:thm:affine-dimension}，得到第一个等式。

	在第二个情形，超越次数具有塔式加性
	\(\operatorname{trdeg}_kL=\operatorname{trdeg}_kK+\operatorname{trdeg}_K L\)。
	确实，取 \(K/k\) 的超越基和 \(L/K\) 的超越基，它们的并在 \(k\) 上代数独立；而 \(L\) 对它们生成的域代数，因为两段剩余扩张都代数。再将 \(A,B\) 的维数分别写成其分式域的超越次数，便得结论。
\end{proof}

这个等式计算泛纤维，不能直接把其中的 \(K\) 换成任意点的剩余域。有限映射则对全部纤维有统一的零维结论，且保留了点集以外的有限代数。

\begin{proposition}
\label{g5:prop:finite-fiber-algebra}
	设 \(A\rightarrow B\) 为交换含幺环同态，且 \(B\) 是有限 \(A\) 模。对每个 \(\mathfrak p\in\operatorname{Spec}A\)，纤维环 \(B(\mathfrak p)\) 是有限维 \(\kappa(\mathfrak p)\) 代数。若它非零，则为零维 Artin 环，只有有限多个素理想，且各剩余域都是 \(\kappa(\mathfrak p)\) 的有限扩张。若 \(A\subseteq B\) 单射，则每个纤维环均非零。
\end{proposition}
\begin{proof}
	取 \(B\) 的有限组 \(A\) 模生成元。张量 \(\kappa(\mathfrak p)\) 后，它们的像张成 \(B(\mathfrak p)\)，所以这是有限维向量空间。它的每个理想都是向量子空间，故理想降链终止。\cref{g5:thm:artinian-zero-dimension}说明非零纤维环维数为零且素理想只有有限多个；每个剩余域还是原有限维向量空间的商，因此为有限域扩张。

	若 \(A\subseteq B\) 单射，模有限性使扩张整。由\cref{g5:thm:lying-over}，存在收缩为 \(\mathfrak p\) 的素理想 \(\mathfrak q\subset B\)。商环与局部化的素理想对应使它给出 \(B(\mathfrak p)\) 的一个素理想，故纤维环不能是零环。
\end{proof}

这把\cref{g5:prop:finite-affine-closed-fibres}的闭点结论推广到了任意素点。例如 \(k[t]\rightarrow k[z]\)、\(t\mapsto z^2\) 是秩二的有限自由扩张。在 \(t=0\) 处的纤维环是 \(k[z]/(z^2)\)，只有一个点、维数为零，但长度为二。有限性限制纤维的维数和点数，并不要求纤维约化。

\begin{example}
\label{g5:exm:fiber-dimension-jump}
	设 \(k\) 为代数闭域，考虑
	\[
		\pi:\mathbb A_k^2\longrightarrow\mathbb A_k^2,\qquad
		(x,y)\longmapsto(u,v)=(x,xy).
	\]
	坐标同态 \(A=k[u,v]\rightarrow B=k[x,y]\) 单射，因为不同单项式 \(u^iv^j\) 的像 \(x^{i+j}y^j\) 仍有不同指数。分式域却相同，因为 \(y=v/u\)，所以泛纤维环是 \(k(u,v)\)，维数为零。

	在闭点 \((a,b)\) 处，纤维环为
	\[
		k[x,y]/(x-a,xy-b).
	\]
	代入 \(x=a\) 后，第二个方程成为 \(ay=b\)，得到\cref{g5:tab:fiber-dimension-jump}中的三种情形。
\end{example}

\begin{table}[htbp]
	\centering
	\caption{映射 \(\pi(x,y)=(x,xy)\) 的闭点纤维}
	\label{g5:tab:fiber-dimension-jump}
	\begin{tabular}{lll}
		\toprule
		基点条件 & 纤维环 & 纤维的形状与维数 \\
		\midrule
		\(a\ne0\) & \(k\) & 一个点，维数 \(0\) \\
		\(a=0,\ b=0\) & \(k[y]\) & 一条直线，维数 \(1\) \\
		\(a=0,\ b\ne0\) & 零环 & 空纤维 \\
		\bottomrule
	\end{tabular}
\end{table}

在 \(u\ne0\) 的开集上，逆公式为 \(x=u,y=v/u\)；而整条直线 \(x=0\) 被送到原点。这里 \(A,B\) 都是维数二的整环，泛纤维维数为零，指定点的纤维却可能是一维或为空。这说明维数差公式的泛纤维结论与逐点结论必须区分。

\ProseParagraphBreak

还可以直接检测这个例子不平坦。取理想 \(J=(u,v)\subset A\)。它有呈示
\[
	A\xrightarrow{\,(-v,u)\,}A^2\longrightarrow J\longrightarrow0:
\]
若 \(au+bv=0\)，模去 \(u\) 后由 \(k[v]\) 为整环得 \(u\mid b\)，写 \(b=uc\) 后便有 \(a=-vc\)。张量 \(B\) 得到
\[
	J\otimes_A B\cong B^2/B(-xy,x).
\]
由包含 \(J\hookrightarrow A\) 诱导的映射把 \((a,b)\) 的类送到 \(ax+bxy\)。类 \((-y,1)\) 在核中，但它非零：若它等于 \(c(-xy,x)\)，第二坐标就要求 \(1=cx\)，这在 \(k[x,y]\) 中不可能。因此 \(J\otimes_A B\rightarrow B\) 不是单射，\(B\) 不是平坦 \(A\) 模。这个判断使用了张量关系，未把纤维维数变化本身当作平坦性判据。

\ProseParagraphBreak

参数系还给出一个适用于指定点的局部维数估计。它的等号条件可直接用素链能否向下提升来表达，因而与\secref{g5:sec:0301}的 going-down 性质相接。

\begin{proposition}
\label{g5:prop:local-fiber-dimension-inequality}
	设 \((A,\mathfrak m)\rightarrow(B,\mathfrak n)\) 为非零 Noether 局部环之间的局部同态，即 \(\mathfrak mB\subseteq\mathfrak n\)。则
	\[
		\dim B\leq\dim A+\dim(B/\mathfrak mB).
	\]
	若该同态具有 going-down 性质，即对 \(\mathfrak p_1\subseteq\mathfrak p_2\) 及收缩为 \(\mathfrak p_2\) 的 \(\mathfrak q_2\in\operatorname{Spec}B\)，总有 \(\mathfrak q_1\subseteq\mathfrak q_2\) 收缩为 \(\mathfrak p_1\)，则上式取等号。
\end{proposition}
\begin{proof}
	记 \(d=\dim A\)、\(r=\dim(B/\mathfrak mB)\)。选 \(A\) 的参数系 \(x_1,\ldots,x_d\)，再选 \(B/\mathfrak mB\) 的参数系，并将其提升为 \(y_1,\ldots,y_r\in\mathfrak n\)。因 \(\sqrt{(x_1,\ldots,x_d)A}=\mathfrak m\)，每个 \(\mathfrak m\) 中元素的某个幂属于 \((x_1,\ldots,x_d)\)。因此
	\[
		\mathfrak mB\subseteq\sqrt{(x_1,\ldots,x_d)B}.
	\]
	而商环中参数系的条件给出
	\(\sqrt{\mathfrak mB+(y_1,\ldots,y_r)B}=\mathfrak n\)。
	两者结合说明 \(x_i,y_j\) 这 \(d+r\) 个元素在 \(B\) 中生成的理想根式为 \(\mathfrak n\)。由\cref{g5:thm:system-of-parameters}的最小性，得到 \(\dim B\leq d+r\)。

	若 going-down 成立，取 \(B/\mathfrak mB\) 中长度为 \(r\) 的素链，提升为
	\[
		\mathfrak q_0\subsetneq\cdots\subsetneq\mathfrak q_r=\mathfrak n.
	\]
	这些素理想都收缩为 \(\mathfrak m\)。再取 \(A\) 中长度为 \(d\)、终点为 \(\mathfrak m\) 的素链，从 \(\mathfrak q_0\) 开始反复使用 going-down，便在 \(\mathfrak q_0\) 下方添上 \(d\) 次严格包含；严格性由收缩理想严格不同保证。这样得到 \(B\) 中长度为 \(d+r\) 的素链，给出相反不等式。
\end{proof}

这正是\cite[第15节，定理15.1]{Matsumura1989}中局部维数估计及 going-down 等号的证明内容。对于一般有限型映射，哪些开集上纤维维数稳定、平坦性怎样控制维数及重数，还需要相对结构定理；本节已经把这些问题分成可计算的泛纤维环、指定点的纤维环及其局部关系。

% 纲要条目排为编号小节标题，并进入目录与书签；说明段仍保留于纲要注释中。


\begin{chaptersummary}
	元素整于基环，等价于它生成的代数在基环上为有限模。这个判据使有限个整元能够共同处理，并导出整性的传递性。整扩张的素理想满足 lying-over、going-up 与 incomparability；整环扩张在基环整闭时还满足 going-down。上行和下行的出发数据不同，这一区别决定能否控制某个指定素理想下方的链。

	\ProseParagraphBreak

	Noether 正规化在有限型代数内部选取代数无关的参数，使原代数成为所选多项式子代数上的有限模。无限域上可以使用线性代换，有限域上则可用彼此分离的幂次使关系首一化。普通正规化是在给定分式域中补入整元；两种构造的方向与目的各不相同。

	\ProseParagraphBreak

	Krull 维数计算素链的长度，素理想的高度等于相应局部环的维数。主理想定理把一个方程上极小素理想的高度限制为一，广义高度定理把这个界推广到方程数。有限型整代数的维数等于分式域的超越次数；一般有限型代数取各极小分支的最大值。整扩张保持整体维数，而指定素理想的高度比较还要核对下降性质。

	\ProseParagraphBreak

	Noether 局部环的参数系用恰好 \(\dim R\) 个元素切出有限长度的商环。嵌入维数计算极大理想所需的生成元数，商环长度记录切片的代数厚度，两者均不能与参数数目混同。有限映射的非空纤维为零维有限代数；一般有限型映射的泛纤维维数由超越次数给出，特殊点上的纤维则可能改变维数、点数或幂零结构。
\end{chaptersummary}

\BookExercises

\begin{exercise}\label{g5:ex:integral-two-square-roots}
	设 \(k\) 为域，\(A=k[u,v]\)，
	\[
		B=A[s,t]/(s^2-u,t^2-v).
	\]
	证明 \(B\) 是秩为四的自由 \(A\) 模，并为 \(z=s+t\) 写出一个首一整性方程。说明该方程在特征二时仍成立，但不能据此断言它是最小次数的方程。
\end{exercise}
\begin{solution}
	先模去 \(s^2-u\)，再模去 \(t^2-v\)，两次首一多项式除法都给出唯一余式。因此每个元素唯一写成
	\[
		a+bs+ct+dst,\qquad a,b,c,d\in A,
	\]
	故 \(1,s,t,st\) 是一组自由基。由
	\[
		z^2-u-v=2st
	\]
	得
	\[
		z^4-2(u+v)z^2+(u-v)^2=0.
	\]
	这是对所有特征都成立的整系数恒等式，所以在特征二时依然有效。但在特征二时已经有 \(z^2=u+v\)，上面的四次式只是这个二次式的平方。有限模判据产生的首一方程保证存在性，不保证次数极小。
\end{solution}

\begin{exercise}\label{g5:ex:integral-localization-units}
	设 \(A\) 为整环，\(S\subseteq A\setminus\{0\}\) 为乘法集。证明下列条件等价：\(S^{-1}A\) 整于 \(A\)；\(S^{-1}A\) 是有限 \(A\) 模；\(S\) 中每个元素都已是 \(A\) 的单位；自然同态 \(A\to S^{-1}A\) 为同构。
\end{exercise}
\begin{solution}
	有限模蕴含整性。若局部化整于 \(A\)，对任意 \(s\in S\)，存在
	\[
		(1/s)^n+a_1(1/s)^{n-1}+\cdots+a_n=0,\qquad a_i\in A.
	\]
	乘以 \(s^{n-1}\)，得到
	\[
		1/s=-a_1-a_2s-\cdots-a_ns^{n-1}\in A.
	\]
	整环假设使 \(A\) 与局部化都能视为同一分式域中的子环，因此这个等式表明 \(s\) 已在 \(A\) 中可逆。若每个分母都是单位，局部化的逆同态就是 \(a/s\mapsto as^{-1}\)。最后，同构于 \(A\) 的模当然由一个元素生成，完成循环。
\end{solution}

\begin{exercise}\label{g5:ex:integral-base-change}
	设 \(A\to B\) 和 \(A\to C\) 为交换含幺环同态。证明：若 \(B\) 是有限 \(A\) 模，则 \(B\otimes_A C\) 是有限 \(C\) 模；若 \(B\) 整于 \(A\) 的像，则 \(B\otimes_A C\) 整于 \(C\) 的像。再计算
	\(\mathbb Z[i]\otimes_{\mathbb Z}\mathbb F_2\)，说明换基后的代数可能出现幂零元。
\end{exercise}
\begin{solution}
	若 \(b_1,\ldots,b_r\) 生成 \(B\)，则 \(b_j\otimes1\) 生成张量积，因为
	\[
		b\otimes c=\sum_j(b_j\otimes1)c_jc
	\]
	对 \(b=\sum_j a_jb_j\) 成立，这里的 \(c_j\) 是 \(a_j\) 在 \(C\) 中的像。

	在整性情形，每个 \(b\otimes1\) 满足由 \(b\) 的首一方程换系数得到的首一方程。张量积中的任意元素都是有限和 \(\sum_{j=1}^r b_j\otimes c_j\)，属于由这些有限个整元生成的 \(C\) 子代数。该子代数是有限 \(C\) 模，因而其中每个元素都整。这也涵盖结构同态有核的情形：方程系数只取决于 \(C\) 的像。

	最后，
	\[
		\mathbb Z[i]\otimes_{\mathbb Z}\mathbb F_2
		\cong\mathbb F_2[T]/(T^2+1)
		=\mathbb F_2[T]/((T+1)^2)
		\cong\mathbb F_2[\varepsilon]/(\varepsilon^2).
	\]
	\(\varepsilon\) 的类非零而平方为零。因此有限自由扩张在换基后仍有限自由，却未必保持约化性。
\end{solution}

\begin{exercise}\label{g5:ex:cusp-two-normalizations}
	设 \(k\) 为域，\(B=k[t^2,t^3]\subset k[t]\)。分别证明 \(B\) 在 \(k[t^2]\) 和 \(k[t^3]\) 上有限自由，求出秩与一组基，并求 \(B\) 在其分式域中的整闭包。说明这三个环扩张对应哪些不同的正规化构造。
\end{exercise}
\begin{solution}
	\(B\) 中出现的单项式指数为 \(0,2,3,4,\ldots\)。按奇偶分类，偶次幂属于 \(k[t^2]\)，奇次幂从 \(t^3\) 开始，均为 \(t^3\) 乘以 \(t^2\) 的幂。因此
	\[
		B=k[t^2]\oplus k[t^2]t^3.
	\]
	按模三余数分类，指数分别从 \(0,4,2\) 开始，故
	\[
		B=k[t^3]\oplus k[t^3]t^2\oplus k[t^3]t^4.
	\]
	两处的独立性均来自单项式指数集合互不相交，所以秩分别为二和三。

	又有 \(t=t^3/t^2\in\operatorname{Frac}B\)，故 \(\operatorname{Frac}B=k(t)\)。元素 \(t\) 满足 \(X^2-t^2=0\)，整于 \(B\)，却不属于 \(B\)。因此 \(B\) 不整闭。\(k[t]\) 是唯一分解整环，按\cref{g5:prop:ufd-integrally-closed}整闭。若 \(h\in k(t)\) 整于 \(B\)，同一个首一方程也表明它整于 \(k[t]\)，所以 \(h\in k[t]\)。反之，\(k[t]=B[t]\) 整于 \(B\)。于是整闭包恰为 \(k[t]\)。

	前两个子环 \(k[t^2]\)、\(k[t^3]\) 都给出 Noether 正规化，参数数目均为一，但模秩不同；扩张 \(B\subset k[t]\) 则是普通正规化，它在原分式域中补入缺失的整元。
\end{solution}

\begin{exercise}\label{g5:ex:finite-field-rank-six}
	在
	\[
		B=\mathbb F_2[x,y]/(xy(x-y))
	\]
	中令 \(u=x^2+xy+y^2\)。证明 \(u\) 在 \(\mathbb F_2\) 上超越，且 \(B\) 是 \(\mathbb F_2[u]\) 上秩为六的自由模，给出一组基。将其与\cref{g5:exm:finite-field-nonlinear-normalization}的秩七投影比较。
\end{exercise}
\begin{solution}
	引入独立变量 \(U,X,Y\)，令
	\[
		g_1=Y^2+XY+X^2+U,\qquad f=X^2Y+XY^2.
	\]
	则
	\[
		B\cong\mathbb F_2[U,X,Y]/(g_1,f).
	\]
	在特征二下，\(f+Xg_1=X^3+UX\)，故
	\[
		(g_1,f)=(g_1,X^3+UX).
	\]
	先对首一三次式 \(X^3+UX\) 取商，再对关于 \(Y\) 的首一二次式 \(g_1\) 取商，得到
	\[
		B\cong
		\bigl(\mathbb F_2[U][X]/(X^3+UX)\bigr)[Y]/
		(Y^2+XY+X^2+U).
	\]
	两次首一除法说明 \(B\) 是 \(\mathbb F_2[U]\) 上的自由模，基为
	\[
		1,\ x,\ x^2,\ y,\ xy,\ x^2y.
	\]
	特别地，\(\mathbb F_2[U]\to B\) 单射，其像即 \(\mathbb F_2[u]\)，故 \(u\) 超越。这里与正文的参数 \(x-y^3\) 都只需一个代数无关元，却得到不同的自由秩。Noether 正规化确定参数数目，并不把所选有限投影的次数固定下来。
\end{solution}

\begin{exercise}\label{g5:ex:mixed-codimension-equations}
	设 \(k\) 为域，\(r\geq2\)，\(S=k[x_0,x_1,\ldots,x_r]\)，
	\[
		I=(x_0x_1,\ldots,x_0x_r).
	\]
	求 \(I\) 的极小素理想及其高度，计算 \(\dim(S/I)\)。再证明，在原点处的局部环中，这 \(r\) 个方程仍是 \(I\) 的极小生成元。
\end{exercise}
\begin{solution}
	有
	\[
		I=(x_0)\cap(x_1,\ldots,x_r).
	\]
	从左到右显然；若 \(x_0g\in(x_1,\ldots,x_r)\)，在商环 \(k[x_0]\) 中得到 \(x_0\overline g=0\)，故 \(g\in(x_1,\ldots,x_r)\)，证明反向包含。这两个素理想互不包含，因而正是 \(I\) 上的极小素理想。

	素链 \((0)\subsetneq(x_0)\) 与主理想定理给出 \(\operatorname{ht}(x_0)=1\)。坐标素链
	\[
		(0)\subsetneq(x_1)\subsetneq\cdots\subsetneq(x_1,\ldots,x_r)
	\]
	给出另一个高度至少为 \(r\)，高度定理给出至多为 \(r\)。两个分支的商环分别为 \(k[x_1,\ldots,x_r]\) 与 \(k[x_0]\)，所以
	\[
		\dim(S/I)=\max\{r,1\}=r.
	\]

	令 \(\mathfrak m=(x_0,\ldots,x_r)\)。在 \(I/\mathfrak mI\) 中，\(x_0x_1,\ldots,x_0x_r\) 的像线性无关，因为 \(\mathfrak mI\) 的每个非零齐次项次数至少为三，而这些二次单项式互不相同。它们已经张成这个商。局部化后剩余向量空间不变，第二章的极小生成元判据便说明至少需要 \(r\) 个生成元。由 \(r\) 个必要方程定义的集合仍然可以同时含余维一与余维 \(r\) 的分支。
\end{solution}

\begin{exercise}\label{g5:ex:parameter-zero-divisor}
	设 \(k\) 为域，
	\[
		R=\bigl(k[x,y]/(x^2,xy)\bigr)_{(x,y)},\qquad
		\mathfrak m=(x,y)R.
	\]
	计算 \(\dim R\) 与嵌入维数，证明 \(y\) 是一个参数，却是零因子，并计算 \(\ell_R(R/yR)\)。
\end{exercise}
\begin{solution}
	\(x\) 幂零，而商去 \((x)\) 得到约化整环 \(k[y]_{(y)}\)，故 \(R\) 的幂零根恰为 \((x)\)。素链在商去幂零根时不变，所以 \(\dim R=1\)。\(\mathfrak m/\mathfrak m^2\) 由 \(x,y\) 的像组成一组 \(k\) 基：原关系 \(x^2,xy\) 没有一次项，局部化不会增加一次线性关系。因此嵌入维数为二。

	又有
	\[
		R/yR\cong k[x]/(x^2),
	\]
	其唯一素理想是 \((x)\)，所以 \(\sqrt{yR}=\mathfrak m\)，\(y\) 构成一个参数系。这里 \(x\ne0\)：它在 \(k[x,y]/(x^2,xy)\) 中的零化子为 \((x,y)\)，故原点之外的分母不能把它杀掉。但 \(yx=0\)，所以 \(y\) 是零因子。商环有组成列
	\[
		0\subset (x)\subset k[x]/(x^2),
	\]
	两个相邻商均同构于 \(k\)，因而长度为二。参数系切出零维局部商，不自动成为正则序列。
\end{solution}

\begin{exercise}\label{g5:ex:artinian-four-dimensional}
	设 \(k\) 为域，\(E=k[x,y]/(x^3,xy,y^2)\)。求其全部素理想、Krull 维数、嵌入维数与长度，并写出极大理想的幂滤过。比较这些数各自测量的对象。
\end{exercise}
\begin{solution}
	首一单项式关系给出 \(k\) 基 \(1,x,x^2,y\)。\(x,y\) 都幂零，所以每个素理想都包含它们；又 \(E/(x,y)\cong k\)，故唯一素理想也是唯一极大理想
	\(\mathfrak m=(x,y)\)。于是 \(E\) 为零维局部环，且
	\[
		E\supset\mathfrak m\supset\mathfrak m^2=(x^2)
		\supset\mathfrak m^3=0.
	\]
	各层 \(E/\mathfrak m,\mathfrak m/\mathfrak m^2,\mathfrak m^2\) 的 \(k\) 维数依次为 \(1,2,1\)。因此嵌入维数为二，长度为 \(1+2+1=4\)。Krull 维数只看素链，此处为零；嵌入维数记录极大理想的极小生成元数；长度还累计了幂零方向保留的全部组成因子。
\end{solution}

\begin{exercise}\label{g5:ex:parameter-ideal-length-six}
	设 \(k\) 为域，\(R=k[x,y]_{(x,y)}\)，\(\mathfrak m=(x,y)R\)，\(J=(x^2,y^3)R\)。证明 \(x^2,y^3\) 为参数系，而 \(J\ne\mathfrak m\)。计算 \(\dim R\)、嵌入维数以及 \(\ell_R(R/J)\)。
\end{exercise}
\begin{solution}
	素链 \((0)\subsetneq(x)R\subsetneq\mathfrak m\) 给出维数至少为二，高度定理给出 \(\dim R\leq2\)，故维数为二。\(x,y\) 的像组成 \(\mathfrak m/\mathfrak m^2\) 的基，所以嵌入维数也是二。由 \(x^2,y^3\in J\)，有 \(\sqrt J=\mathfrak m\)，故两个元素构成参数系。但 \(x\notin J\)，所以参数理想不等于极大理想。

	多项式商环 \(k[x,y]/(x^2,y^3)\) 有基
	\[
		1,\ y,\ y^2,\ x,\ xy,\ xy^2.
	\]
	其中极大理想 \((x,y)\) 幂零，任何常数项非零的元素 \(c+n\) 都可逆：将其写成 \(c(1+c^{-1}n)\)，逆由有限几何级数给出。因此原点处的局部化不再改变这个商环，\(R/J\) 的 \(k\) 维数为六。它的剩余域为 \(k\)，故 \(\ell_R(R/J)=6\)。参数个数为二，切片长度为六，二者记录不同的信息。
\end{solution}

\begin{exercise}\label{g5:ex:finite-local-base-maximals}
	设 \((A,\mathfrak m)\) 为整闭 Noether 局部整环，\(A\subseteq B\) 为有限扩张，且 \(B\) 为整环。证明 \(B\) 只有有限多个极大理想，每个极大理想 \(\mathfrak n\) 都满足
	\[
		\dim B_{\mathfrak n}=\dim A.
	\]
	这是否意味着 \(B\) 本身为局部环？以
	\[
		A=\mathbb Z_{(5)},\qquad B=A[T]/(T^2+1)
	\]
	为例，求 \(B\) 的全部极大理想和相应局部维数。
\end{exercise}
\begin{solution}
	有限性蕴含整性，由\cref{g5:lem:integral-field-extension}，\(B\) 的每个极大理想都收缩为 \(A\) 的极大理想 \(\mathfrak m\)，因此都包含 \(\mathfrak mB\)。这些极大理想恰与有限维剩余域代数
	\[
		B/\mathfrak mB=B\otimes_A(A/\mathfrak m)
	\]
	的极大理想对应。由\cref{g5:prop:finite-fiber-algebra}，它们只有有限多个。又因 \(A\) 整闭且 \(B\) 为整环，going-down 成立；\cref{g5:prop:integral-height-comparison}于是给出
	\[
		\dim B_{\mathfrak n}
		=\operatorname{ht}_B\mathfrak n
		=\operatorname{ht}_A\mathfrak m
		=\dim A.
	\]
	这里既用有限性控制极大理想个数，又用下降性质控制每个指定极大理想的高度。

	在所给例子中，\(A\) 为主理想整环的局部化，故 Noether 且整闭。\(B\) 是秩二的自由 \(A\) 模，并可识别为 \(\mathbb Z[i]\) 在所有不被 \(5\) 整除的整数处的局部化，因此也是整环。模去 \(5\) 得
	\[
		B/5B\cong\mathbb F_5[T]/((T-2)(T+2))
		\cong\mathbb F_5\times\mathbb F_5.
	\]
	故全部极大理想恰为
	\[
		\mathfrak n_+=(5,T-2),\qquad
		\mathfrak n_-=(5,T+2).
	\]
	两个局部环的维数都为一，但 \(B\) 有两个极大理想，因而不是局部环。
\end{solution}

\begin{exercise}\label{g5:ex:finite-extra-point}
	设 \(k\) 为域，令 \(A=k[t]\)，\(B=A\times k\)，结构同态为
	\[
		A\longrightarrow B,\qquad a(t)\longmapsto(a(t),a(0)).
	\]
	证明这是有限单射，计算 \(\dim A,\dim B\)。令 \(\mathfrak q=A\times(0)\)，计算其收缩与高度。再计算各 \(a\in k\) 处的纤维环 \(B\otimes_A A/(t-a)\)，并证明 \(B\) 不是平坦 \(A\) 模。
\end{exercise}
\begin{solution}
	第一分量保证单射；元素 \((1,0),(0,1)\) 生成 \(B\) 作为 \(A\) 模，因为
	\[
		(f,c)=f\cdot(1,0)+c\cdot(0,1).
	\]
	因此扩张有限且整，\(\dim B=\dim A=1\)。

	\(\mathfrak q\) 的商环为 \(k\)，收缩是 \((t)\)。乘积环中的素理想必含两个正交幂等元 \((1,0),(0,1)\) 之一，因为它们的乘积为零；于是素理想只有形如 \(\mathfrak p\times k\) 或 \(A\times(0)\) 的两类。因此 \(\mathfrak q\) 下方没有更小的素理想，\(\operatorname{ht}\mathfrak q=0\)，而 \(\operatorname{ht}(t)=1\)。这不违背整体维数保持；前题所需的上环为整环这一条件在这里失败。

	纤维环就是 \(B/(t-a)B\)。在第一分量取商得到 \(k\)；第二分量上 \(t-a\) 作用为标量 \(-a\)。故
	\[
		B\otimes_A A/(t-a)\cong
		\begin{cases}
			k,&a\ne0,\\
			k\times k,&a=0.
		\end{cases}
	\]
	最后，\(A\xrightarrow{\times t}A\) 单射，张量 \(B\) 后变为 \(B\) 上乘以 \((t,0)\)。它杀掉非零元素 \((0,1)\)，所以不再单射，\(B\) 不平坦。
\end{solution}

\begin{exercise}\label{g5:ex:gaussian-prime-fibers}
	把 \(\mathbb Z[i]\) 视为 \(\mathbb Z\) 上的有限代数。分别计算素数 \(p=2,3,5\) 处的纤维环，求纤维中素理想的个数及其剩余域，并比较纤维的维数、\(\mathbb F_p\) 向量空间维数与约化性。
\end{exercise}
\begin{solution}
	\(\mathbb Z[i]\) 在 \(\mathbb Z\) 上有基 \(1,i\)，故每个纤维都有 \(\mathbb F_p\) 维数二，且
	\[
		\mathbb Z[i]\otimes_{\mathbb Z}\mathbb F_p
		\cong\mathbb F_p[T]/(T^2+1).
	\]
	当 \(p=2\) 时，\(T^2+1=(T+1)^2\)，纤维为
	\(\mathbb F_2[\varepsilon]/(\varepsilon^2)\)。它只有一个素理想 \((\varepsilon)\)，剩余域为 \(\mathbb F_2\)，但不是约化环。

	当 \(p=3\) 时，\(\mathbb F_3\) 中的平方只有 \(0,1\)，因此 \(T^2+1\) 无根且不可约。纤维是九元域 \(\mathbb F_9\)，只有零素理想，剩余域为 \(\mathbb F_9\)。

	当 \(p=5\) 时，\(T^2+1=(T-2)(T+2)\)，两个因子互素，中国剩余定理给出 \(\mathbb F_5\times\mathbb F_5\)。纤维有两个素理想，剩余域均为 \(\mathbb F_5\)。

	三个纤维的 Krull 维数都为零，向量空间维数都为二，素理想个数却依次为一、一、二；只有第一个不约化。顺便由整扩张的维数保持得 \(\dim\mathbb Z[i]=\dim\mathbb Z=1\)，其中 \(\mathbb Z\) 的非零素理想均为极大理想。
\end{solution}

\begin{exercise}\label{g5:ex:flat-hyperbola-family}
	设 \(k\) 为域，\(A=k[t]\)，\(B=k[x,y]\)，结构同态由 \(t\mapsto xy\) 给出。证明 \(B\) 是自由 \(A\) 模，且一组基为
	\[
		1,\quad x^i\ (i\geq1),\quad y^j\ (j\geq1).
	\]
	由此证明 \(B\) 平坦但不是有限 \(A\) 模。计算泛纤维及每个 \(t=a\)、\(a\in k\) 处的纤维，并说明平坦性与特殊纤维可约并不冲突。
\end{exercise}
\begin{solution}
	不同的 \(t\) 次幂映到不同的单项式 \(x^ny^n\)，所以结构同态单射。每个单项式 \(x^ay^b\) 都可以取出 \(t^{\min(a,b)}\)，剩下 \(1\)、正次 \(x\) 幂或正次 \(y\) 幂，证明所列元素张成 \(B\)。对于独立性，展开系数多项式后，三个族的单项式指数分别位于 \(a=b\)、\(a>b\)、\(a<b\)；在每一族内，指数又唯一确定系数中的 \(t\) 次幂及基元素。因此全部单项式互不重复，任何有限线性关系的系数都为零。

	张量这个自由模等于取相应指标集上的直和。直和逐坐标保留单射，故 \(B\) 平坦。无限基的自由模不能有限生成：有限组元素只涉及有限多个基方向，不能生成剩余方向。因此这个扩张不有限。

	令 \(K=k(t)\)。因 \(t\) 在 \(K\) 中非零且可逆，
	\[
		B\otimes_AK\cong K[x,y]/(xy-t)\cong K[x,x^{-1}],
	\]
	同构由 \(y=t/x\) 给出。对 \(a\ne0\)，闭点纤维同样为
	\[
		k[x,y]/(xy-a)\cong k[x,x^{-1}],
	\]
	而 \(a=0\) 的纤维为 \(k[x,y]/(xy)\)，其两个极小素理想为 \((x)\)、\((y)\)。各环的维数都为一，泛纤维和非零参数纤维为整环，零参数纤维则有两个分支。这里平坦性由自由模结构直接证明，并不要求每个纤维都不可约。
\end{solution}

\begin{exercise}\label{g5:ex:zero-dimensional-finite-type}
	设 \(B\ne0\) 为域 \(k\) 上的有限型代数。证明
	\[
		\dim B=0\quad\Longleftrightarrow\quad\dim_k B<\infty.
	\]
	说明允许 \(B\) 含非零幂零元，并用 \(k(t)\) 解释有限型假设为何必要。
\end{exercise}
\begin{solution}
	由 Noether 正规化，存在多项式子代数 \(C=k[u_1,\ldots,u_d]\subseteq B\)，使 \(B\) 为有限 \(C\) 模。整扩张保持维数，故
	\[
		\dim B=\dim C=d.
	\]
	若 \(\dim B=0\)，则 \(d=0\)，\(C=k\)，所以 \(B\) 为有限维 \(k\) 向量空间。反之，若 \(B\) 有限维，全部理想都是向量子空间，因此降链终止，\(B\) 为 Artin 环，由\cref{g5:thm:artinian-zero-dimension}得 \(\dim B=0\)。两向证明均未使用约化性；例如 \(k[\varepsilon]/(\varepsilon^n)\)、\(n\geq2\)，满足等价条件。

	域 \(k(t)\) 的 Krull 维数为零，但 \(1,t,t^2,\ldots\) 在 \(k\) 上线性无关，故其 \(k\) 维数无穷。由刚证的等价关系可知，它不能是有限型 \(k\) 代数。环作为模有限与作为代数有限型的区别，在这个例子中均不可略去。
\end{solution}

\begin{exercise}\label{g5:ex:nonnoetherian-principal-height}
	设 \(k\) 为域，
	\[
		A_0=k[x]_{(x)},\qquad K=k(x),\qquad V=K[y]_{(y)}.
	\]
	记 \(\varepsilon:V\to K\) 为 \(y=0\) 处的取值同态，定义
	\[
		R=\varepsilon^{-1}(A_0),\qquad
		\mathfrak p=yV,\qquad
		\mathfrak m=\varepsilon^{-1}(xA_0)\subset R.
	\]
	证明 \(R\) 为局部整环，\(\mathfrak m=xR\) 为其极大理想，且
	\[
		(0)\subsetneq\mathfrak p\subsetneq\mathfrak m
	\]
	为素链。再构造 \(R\) 中严格上升的理想链，说明主理想定理的 Noether 假设不能删去。
\end{exercise}
\begin{solution}
	\(V\) 是整环，故其子环 \(R\) 也是整环。映射 \(\varepsilon:R\to A_0\) 满射，因为每个 \(A_0\) 元素都可以视为与 \(y\) 无关的常数。其核为 \(yV=\mathfrak p\)，而 \(\mathfrak p\subset R\)，所以
	\[
		R/\mathfrak p\cong A_0,\qquad R/\mathfrak m\cong k.
	\]
	因此 \(\mathfrak p\) 为素理想，\(\mathfrak m\) 为极大理想。

	若 \(f\in R\setminus\mathfrak m\)，则 \(\varepsilon(f)\) 为 \(A_0\) 的单位，特别地非零。这使 \(f\) 在 \(V=K[y]_{(y)}\) 中可逆，且
	\[
		\varepsilon(f^{-1})=\varepsilon(f)^{-1}\in A_0,
	\]
	故 \(f^{-1}\in R\)。反之，\(R\) 的单位在 \(A_0\) 中的像也必须可逆。因此全部非单位恰为 \(\mathfrak m\)，\(R\) 是局部环。

	显然 \(xR\subseteq\mathfrak m\)。若 \(f\in\mathfrak m\)，则 \(f/x\in V\)，因为 \(x\) 是 \(K\) 中的非零常数；又
	\(\varepsilon(f/x)=\varepsilon(f)/x\in A_0\)，故 \(f/x\in R\)。这给出 \(\mathfrak m=xR\)。元素 \(y\ne0\) 属于 \(\mathfrak p\)，元素 \(x\) 属于 \(\mathfrak m\setminus\mathfrak p\)，所列素链严格。因此一个主理想本身就是高度至少为二的极大素理想。

	最后，每个 \(y/x^n\) 都属于 \(R\)，因为其取值为零。于是有上升链
	\[
		(y)R\subsetneq(y/x)R\subsetneq(y/x^2)R\subsetneq\cdots.
	\]
	包含关系来自 \(y/x^n=x(y/x^{n+1})\)。若某一步相等，则 \(y/x^{n+1}=r(y/x^n)\) 对某个 \(r\in R\) 成立。在整环的分式域中约去非零因子，得到 \(r=1/x\)，但 \(\varepsilon(1/x)=1/x\notin A_0\)，矛盾。故 \(R\) 不为 Noether 环，确实处在高度定理的假设之外。
\end{solution}
